% \iffalse meta-comment
% Copyright (C) 2026 Andreas Willibald Weber
%
% This work may be distributed and/or modified under the conditions of the
% LaTeX Project Public License, either version 1.3c of this license or (at your
% option) any later version. The latest version of the license is available at
% https://www.latex-project.org/lppl.txt.
%
% This work has the LPPL maintenance status `maintained'. The Current
% Maintainer of this work is Andreas Willibald Weber.
%
% This work consists of the source files kanban.dtx and kanban.ins; build.lua,
% README.md, CONTRIBUTING.md, RELEASING.md, CHANGELOG.md, and LICENSE; the
% example files kanban-example.tex, kanban-modes-example.tex,
% kanban-deadlines-example.tex, kanban-date-formats-example.tex,
% kanban-palette-example.tex, kanban-columnsep-example.tex,
% kanban-break-example.tex, and kanban-landscape-example.tex; the files in the
% testfiles directory; and the derived files kanban.sty and kanban.pdf.
% \fi
%
% \iffalse
%<*driver>
\documentclass[a4paper,11pt]{ltxdoc}
\makeatletter
\let\kanban@docverbatim\verbatim
\let\kanban@enddocverbatim\endverbatim
\let\kanban@doc@verbatim\@verbatim
\let\kanban@doc@xverbatim\@xverbatim
\makeatother
\usepackage[T1]{fontenc}
\usepackage{lmodern}
\usepackage[
  left=40mm,
  right=32mm,
  top=34mm,
  bottom=33mm,
  headheight=15pt,
  marginparwidth=30mm,
  marginparsep=4mm
]{geometry}
\usepackage{microtype}
\usepackage{ragged2e}
\usepackage{enumitem}
\usepackage{fancyhdr}
\usepackage{titlesec}
\usepackage{needspace}
\usepackage{listings}
\usepackage{pdflscape}
\usepackage{kanban}
\tcbuselibrary{breakable}
\newcommand{\kanbanManualBreakCard}[3]{%
  \kanbantask[author=#2]{#1}{#3. Document the assumptions, record the
    outcome, and retain the evidence required for the next project review.}}
\newcommand{\kanbanManualBreakDeadlineCard}[4]{%
  \kanbantask[author=#2,#3]{#1}{#4. Document the assumptions, record the
    outcome, and retain the evidence required for the next project review.}}
\newcommand{\kanbanManualDateFormatBoard}[2]{%
  \begin{minipage}[t]{0.48\linewidth}
    \textit{#2}\par\smallskip
    \begin{kanban}[columns={Upcoming},date-format=#1]
      \kanbantask[due=2099-07-18]{Upcoming}{Prepare the scheduled release}
    \end{kanban}
  \end{minipage}}
\definecolor{kanbanManualRule}{HTML}{D7DEE2}
\definecolor{kanbanManualCodeBack}{HTML}{F3F4F5}
\definecolor{kanbanManualCodeRule}{HTML}{D8DDE1}
\lstdefinestyle{kanbanmanual}{
  language={[LaTeX]TeX},
  basicstyle=\fontsize{9pt}{11pt}\selectfont\ttfamily\color{kanbanInk},
  keywordstyle=\color{kanbanBlue!85!black},
  commentstyle=\itshape\color{kanbanSage!75!black},
  backgroundcolor=\color{kanbanManualCodeBack},
  frame=single,
  rulecolor=\color{kanbanManualCodeRule},
  framerule=0.35pt,
  framesep=0pt,
  framexleftmargin=2.5mm,
  framexrightmargin=2.5mm,
  framextopmargin=1.2mm,
  framexbottommargin=1.2mm,
  aboveskip=0.8\baselineskip,
  belowskip=0.8\baselineskip,
  columns=fullflexible,
  keepspaces=true,
  gobble=2,
  breaklines=true,
  showstringspaces=false,
  tabsize=2
}
\lstset{style=kanbanmanual}
\BeforeBeginEnvironment{lstlisting}{\Needspace{7\baselineskip}}
\tcolorboxenvironment{macrocode}{
  enhanced,
  breakable,
  sharp corners,
  colback=kanbanManualCodeBack,
  frame hidden,
  boxrule=0pt,
  boxsep=0pt,
  left=0pt,
  right=0pt,
  top=1mm,
  bottom=1mm,
  before skip=0.8\baselineskip,
  after skip=0.8\baselineskip
}
\tcolorboxenvironment{macrocode*}{
  enhanced,
  breakable,
  sharp corners,
  colback=kanbanManualCodeBack,
  frame hidden,
  boxrule=0pt,
  boxsep=0pt,
  left=0pt,
  right=0pt,
  top=1mm,
  bottom=1mm,
  before skip=0.8\baselineskip,
  after skip=0.8\baselineskip
}
\setlength{\parindent}{0pt}
\setlength{\parskip}{0.55em plus 0.1em minus 0.1em}
\raggedbottom
\setlength{\MacrocodeTopsep}{0.8\baselineskip}
\setlist{itemsep=0.3em,topsep=0.45em,parsep=0pt}
\makeatletter
\let\@verbatim\kanban@doc@verbatim
\let\@xverbatim\kanban@doc@xverbatim
\let\verbatim\kanban@docverbatim
\let\endverbatim\kanban@enddocverbatim
\apptocmd{\@verbatim}
  {\advance\leftskip 2.5mm\advance\rightskip 2.5mm}
  {}{\PackageWarning{kanban-doc}{Could not adjust verbatim indentation}}
\renewcommand{\IndexParms}{%
  \raggedcolumns
  \parindent=0pt
  \columnsep=8pt
  \parskip=0pt plus 1pt
  \rightskip=0pt plus 2em
  \parfillskip=0pt
  \emergencystretch=1em
  \mathsurround=0pt
  \fontsize{7.8pt}{9.6pt}\selectfont
  \def\@idxitem{\par\hangindent 24pt}%
  \def\subitem{\@idxitem\hspace*{12pt}}%
  \def\subsubitem{\@idxitem\hspace*{20pt}}%
  \def\indexspace{\par\vspace{8pt plus 2pt minus 2pt}}}
\makeatother
\setcounter{tocdepth}{2}
\titleformat{\section}
  {\Large\bfseries\sffamily\color{kanbanSlate}}
  {\thesection}{0.7em}{}
\titleformat{\subsection}
  {\large\bfseries\sffamily\color{kanbanBlue}}
  {\thesubsection}{0.7em}{}
\titlespacing*{\section}{0pt}{2.2em}{0.7em}
\titlespacing*{\subsection}{0pt}{1.5em}{0.45em}
\pagestyle{fancy}
\fancyhf{}
\fancyhead[L]{\small\sffamily\textbf{kanban}}
\fancyhead[R]{\small\sffamily User manual}
\fancyfoot[C]{\small\thepage}
\renewcommand{\headrulewidth}{0.4pt}
\renewcommand{\footrulewidth}{0pt}
\AtBeginDocument{%
  \hypersetup{
    linkcolor=kanbanBlue,
    urlcolor=kanbanTeal,
    citecolor=kanbanPlum,
    pdfauthor={Andreas Willibald Weber},
    pdftitle={The kanban package},
    pdfsubject={Kanban project boards in LaTeX}
  }
  \renewcommand{\MacroFont}{%
    \fontsize{8.8pt}{10.5pt}\selectfont\ttfamily\color{kanbanInk}}}
\EnableCrossrefs
\CodelineIndex
\RecordChanges
\DoNotIndex{\NeedsTeXFormat,\ProvidesPackage,\RequirePackage}
\DoNotIndex{\definecolor,\newbool,\newcounter,\newlength,\setlength}
\DoNotIndex{\def,\gdef,\edef,\let,\newcommand,\renewcommand}
\DoNotIndex{\NewDocumentCommand,\NewDocumentEnvironment}
\DoNotIndex{\DeclareOption,\ProcessOptions,\PackageError,\PackageWarning}
\DoNotIndex{\par,\noindent,\ignorespaces,\relax}
\DoNotIndex{\ifblank,\ifbool,\ifboolexpr,\ifdefempty,\ifdefstring}
\DoNotIndex{\hfill,\hspace,\vspace}
\DoNotIndex{\textcolor,\bfseries,\scriptsize,\raggedright}
\DoNotIndex{\kanban@validateboardgeometry}
\begin{document}
  % Keep rendered deadline examples reproducible across rebuilds without
  % altering TeX's calendar primitives for the surrounding manual.
  \kbset{reference-date=2026-07-18}
  \DocInput{kanban.dtx}
\end{document}
%</driver>
% \fi
%
% \GetFileInfo{kanban.sty}
% \changes{v1.0.0}{2026/07/18}{Initial public release.}
%
% \title{The \textsf{kanban} package\\
%   \large Kanban project boards in LaTeX}
% \author{Andreas Willibald Weber}
% \date{Version \fileversion, \filedate}
%
% \hypersetup{pageanchor=false}
% \newgeometry{left=32mm,right=37mm,top=34mm,bottom=33mm}
% \begin{titlepage}
% \centering
% \vspace*{0.10\textheight}
% {\color{kanbanBlue}\rule{0.84\textwidth}{1.2pt}\par}
% \vspace{1.4em}
% {\Huge\bfseries\sffamily\color{kanbanSlate} kanban\par}
% \vspace{0.5em}
% {\Large\sffamily\color{kanbanBlue} Kanban project boards in LaTeX\par}
% \vspace{1.4em}
% {\color{kanbanBlue}\rule{0.84\textwidth}{1.1pt}\par}
% \vfill
% \begin{minipage}{0.84\textwidth}
% \centering
% {\large\bfseries\sffamily Abstract\par}
% \vspace{0.8em}
% \small\justifying
% The \textsf{kanban} package typesets compact project boards composed of
% named columns and task cards. A card may carry a priority, an author, a
% deadline, and a completion marker. Deadlines can be expressed either as
% calendar dates or as months relative to a project start date. Two rendering
% layouts are provided: a fixed horizontal board suited to short summaries and
% a page-breakable layout for longer task lists. Named palettes and light or
% dark board modes can be selected at package, document, or board level.
% Deadline chips distinguish overdue, current, and future values. A reference
% date can be fixed for reproducible output; otherwise the compilation date is
% used. Five common numeric date formats are available, with US notation as the
% initial default.
% This manual describes the public interface, gives complete usage examples,
% and explains the date and configuration models.
% \end{minipage}
% \vfill
% {\large Andreas Willibald Weber\par}
% \vspace{0.4em}
% {\sffamily Version \fileversion\quad--\quad\filedate\par}
% \vspace{2em}
% \end{titlepage}
% \restoregeometry
% \hypersetup{pageanchor=true}
% \pagenumbering{roman}
%
% \tableofcontents
% \clearpage
% \pagenumbering{arabic}
%
% \section{Introduction}
%
% A Kanban board represents work by assigning cards to named stages such as
% ``Backlog'', ``In progress'', and ``Done''. The package follows that model
% directly. A board uses either the built-in columns, the current document
% default, or a board-local list selected with the |columns| environment
% option. The author then assigns each task to exactly one of those columns.
% Columns are rendered in declaration order, and tasks retain their order
% within each column.
%
% The package is intended for reports, project notes, teaching material, and
% other documents in which a board should be generated from ordinary LaTeX
% source. It is not an interactive project-management system: tasks are local
% document data, and the package neither persists nor automatically moves them.
%
% \begin{tcolorbox}[
%   enhanced,
%   title={At a glance},
%   colback=kanbanSlate!4,
%   colframe=kanbanSlate!55,
%   colbacktitle=kanbanSlate,
%   coltitle=white,
%   fonttitle=\bfseries\sffamily,
%   boxrule=0.5pt,
%   arc=1mm,
%   left=4mm,right=4mm,top=3mm,bottom=3mm
% ]
% \begin{description}
% \item[Core environment] |\begin{kanban}[options]|, with board-local columns
%   selected by the |columns={...}| option when required.
% \item[Task command] |\kanbantask[options]{column}{text}|, with |\kbtask|
%   as a concise alias when that short name is available.
% \item[Shared setup] |\kanbansetup{key=value,...}|, with |\kbset| as a
%   concise alias when that short name is available.
% \item[Visual choices] Named palettes, light or dark mode, and fixed or
%   page-breakable layouts.
% \item[Date choices] US, English, German, ISO, and slash-separated year-first
%   display formats, with validated Gregorian classification.
% \item[Column geometry] Synchronized headings and transparent column bodies
%   in the fixed layout, with configurable separation.
% \end{description}
% \end{tcolorbox}
%
% \section{Requirements and installation}
%
% The package requires a LaTeX format dated 2020-10-01 or later and loads
% \textsf{xcolor}, \textsf{etoolbox}, \textsf{tikz}, \textsf{tcolorbox}, and
% \textsf{paracol}. These dependencies are included in the principal TeX
% distributions. Continuous integration runs the regression and documentation
% builds on both a current TeX Live installation and an actual TeX Live 2020
% installation. The declared format boundary is also exercised locally through
% the official \textsf{latexrelease} compatibility mechanism; the regression
% suite is run with pdfTeX, LuaTeX, and XeTeX.
%
% To install from the documented source, run
% \begin{lstlisting}
% latex kanban.ins
% \end{lstlisting}
% and copy |kanban.sty| to a directory searched by TeX. For a private
% installation, that directory is normally a suitable location below
% |TEXMFHOME|. No manual installation is necessary when the package is
% installed through a TeX distribution.
%
% Load the package in the preamble. No options are required:
% \begin{lstlisting}
% \usepackage{kanban}
% \end{lstlisting}
% Initial document defaults may also be set while loading the package:
% \begin{lstlisting}
% \usepackage[
%   columns={Planned,Active,Complete},
%   palette=muted,
%   mode=light,
%   date-format=us,
%   startdate=07/01/2026,
%   reference-date=2026-07-18,
%   columnsep=2mm
% ]{kanban}
% \end{lstlisting}
% These defaults can be changed later with |\kanbansetup| (or its concise
% alias |\kbset|) or overridden by an individual board. The following sections
% explain this hierarchy in detail.
%
% \section{The board environment}
%
% \DescribeEnv{kanban}
% The environment has the syntax
% \begin{quote}
% |\begin{kanban}[|\meta{board options}|]|\\
% \hspace*{2em}\meta{task declarations}\\
% |\end{kanban}|
% \end{quote}
% If |columns| is omitted, or if |columns={}| is supplied, the board uses the
% current document-wide column list. The package initially defines that list as
% |BACKLOG,DOING,DONE|. A non-empty comma-separated value replaces the default
% for this board only. The additional braces around the value protect its
% commas from the surrounding key list. Each name is trimmed and converted to
% uppercase before it is used as a visible heading or a lookup key.
% Task targets undergo the same conversion, so |ColName|, |colname|, |Colname|,
% and |COLNAME| all identify the same column. Column names must be unique after
% normalization and must not contain commas.
%
% Each board is independent. When a board begins, the package first resolves
% the applicable column list. It then creates an empty task list and a zero
% count for every resolved column before collecting task declarations. The
% visible board is produced only when the environment ends. This deferred
% construction is why tasks can be written in any cross-column order while
% still appearing beneath the correct headings.
%
% \subsection{Board options}
%
% The optional argument accepts the following keys.
%
% \begin{description}
% \item[|break|]
%   Selects the page-breakable layout. Each board column becomes a
%   \textsf{paracol} column with an independent colored header followed by its
%   cards. This layout is appropriate when task lists may continue onto a later
%   page. Headers are printed at the start of the board and are not repeated on
%   continuation pages.
% \item[\texttt{columns=\{}\meta{column list}\texttt{\}}]
%   Replaces the current document-wide column list for this board only. The
%   value must be enclosed in braces because it contains commas. An empty value
%   retains the current document default.
% \item[|columnsep=|\meta{length}]
%   Sets the horizontal distance between adjacent columns for this board. The
%   default is |2mm|. This value overrides the document default for the current
%   board only and is used by both the fixed and page-breakable layouts.
% \item[|palette=|\meta{name}]
%   Selects a previously defined color palette for this board. The local
%   selection does not change the palette used by later boards.
% \item[|date-format=us|, |english|, |german|, |iso|, or |ymd|]
%   Selects the numeric format used to print validated calendar dates on this
%   board. The initial default is |us|. This setting also resolves dates whose
%   slash-separated day and month are both at most 12. A local selection does
%   not change the format used by later boards.
% \item[|startdate=|\meta{date}]
%   Sets the project start date for this board. It accepts any supported
%   calendar syntax and is used only when interpreting project-month deadlines
%   such as |M2|. This value overrides the document default for the current
%   board only.
% \item[|reference-date=|\meta{date}]
%   Sets the date against which calendar and project-month deadlines are
%   classified for this board. Without this key, the package uses the TeX
%   compilation date. The exact alias |today| is also accepted. An explicit
%   reference date makes archival and automated builds reproducible.
% \item[|mode=light| or |mode=dark|]
%   Selects the visual mode for this board. The default is |light|. A local
%   mode does not change the mode of subsequent boards.
% \item[\texttt{labels=\{}\meta{label options}\texttt{\}}]
%   Replaces one or more semantic chip labels for this board. The accepted
%   nested keys are |high|, |medium|, |low|, and |done|. Omitted labels retain
%   their current document-wide values.
% \end{description}
%
% The |break| key is board-local by design; it is not accepted as a package
% option or by |\kanbansetup|. Without |break|, the package constructs a
% two-row \textsf{tcolorbox} raster inside one full-width minipage. The first
% row contains the synchronized headings and the second contains the
% synchronized column bodies. The complete board is therefore one horizontal
% object and does not break naturally across pages. This default is preferable
% for compact boards. Task cards occupy the complete available width of the
% column body in both layouts.
%
% \subsection{A first board}
%
% The following complete example applies the environment syntax and default
% fixed layout described above. Because no column list is supplied, it uses
% the built-in |BACKLOG|, |DOING|, and |DONE| stages and assigns one task to
% each stage. The deliberately varied capitalization of the task targets also
% demonstrates case-insensitive assignment.
% \begin{lstlisting}
% \begin{kanban}
%   \kanbantask[prio=high,author=Alice]
%     {backlog}{Agree on the public API}
%   \kanbantask[due=07/31/2026]
%     {DoInG}{Write the implementation}
%   \kanbantask[completed]
%     {DONE}{Create the repository}
% \end{kanban}
% \end{lstlisting}
%
% The example produces the following board:
%
% \begin{kanban}
%   \kanbantask[prio=high,author=Alice]
%     {backlog}{Agree on the public API}
%   \kanbantask[due=07/31/2026]
%     {DoInG}{Write the implementation}
%   \kanbantask[completed]
%     {DONE}{Create the repository}
% \end{kanban}
%
% \section{Configuration hierarchy}
%
% The package uses a consistent hierarchy for column lists, palettes, visual
% modes, date formats, project start dates, reference dates, semantic labels,
% and column separation. This
% makes document-wide settings easy to reuse while keeping local exceptions
% explicit.
%
% \begin{enumerate}
% \item Package options establish the initial defaults.
% \item |\kanbansetup| changes the defaults for boards that follow it.
% \item Environment options override the defaults for one board only, including
%   the board-local |columns| key.
% \end{enumerate}
%
% The most local value has priority. For example:
% \begin{lstlisting}
% \usepackage[
%   columns={Planned,Active,Complete},
%   palette=muted,
%   mode=light,
%   date-format=us,
%   reference-date=2026-07-18,
%   columnsep=2mm
% ]{kanban}
%
% % Uses every package default, including the configured columns.
% \begin{kanban}
%   ...
% \end{kanban}
%
% % Changes defaults for subsequent boards.
% \kbset{
%   columns={Queued,Processing,Stored},
%   palette=monochrome,
%   mode=dark,
%   date-format=iso,
%   reference-date=2026-07-18,
%   columnsep=1.5mm
% }
%
% % Uses the revised document defaults.
% \begin{kanban}
%   ...
% \end{kanban}
%
% % Uses local columns, default palette, light mode, English dates, and a 3 mm gap.
% \begin{kanban}[
%   columns={Todo,Doing,Done},
%   palette=default,mode=light,date-format=english,columnsep=3mm
% ]
%   ...
% \end{kanban}
%
% % Uses Queued, Processing, and Stored again.
% \begin{kanban}
%   ...
% \end{kanban}
% \end{lstlisting}
%
% \DescribeMacro{\kanbansetup}
% The setup command accepts a comma-separated key list:
% \begin{quote}
% |\kanbansetup{columns={|\meta{column list}|},palette=|\meta{name}|,|\\
% \hspace*{2em}|mode=|\meta{mode}|,date-format=|\meta{format}|,|\\
% \hspace*{2em}|startdate=|\meta{date}|,reference-date=|\meta{date}|,|\\
% \hspace*{2em}|columnsep=|\meta{length}|,labels={|\meta{label options}|}}|
% \end{quote}
% When the name is available, the shorter |\kbset| command is an exact alias.
% The longer name is the canonical form because it is less likely to conflict
% with another package. If |\kbset| already exists when \textsf{kanban} is
% loaded, the package preserves the existing definition, issues a warning, and
% leaves |\kanbansetup| fully available.
% A key may be omitted when its current value should remain unchanged. Setup
% calls affect later boards; they do not redraw boards that have already been
% typeset. The |columns| value must be enclosed in braces because it is itself
% a comma-separated list. An empty |columns| value is rejected with a warning,
% and the previous document default remains active.
%
% \subsection{Default column lists}
%
% With no configuration, every board that omits its |columns| option uses the
% three columns |BACKLOG|, |DOING|, and |DONE|. This provides a useful
% conventional board without requiring repeated declarations:
% \begin{lstlisting}
% \begin{kanban}
%   \kanbantask{backlog}{Record the request}
%   \kanbantask{Doing}{Prepare the response}
%   \kanbantask{dOnE}{Archive the decision}
% \end{kanban}
% \end{lstlisting}
%
% A document can establish different defaults when the package is loaded:
% \begin{lstlisting}
% \usepackage[columns={Planned,Active,Complete}]{kanban}
% \end{lstlisting}
% or change them from a particular point onward:
% \begin{lstlisting}
% \kbset{columns={Queued,Processing,Stored}}
% \end{lstlisting}
% In both cases, later environments may omit the |columns| key. Supplying a
% board-local key remains the appropriate way to describe an exceptional board:
% \begin{lstlisting}
% \begin{kanban}[columns={Requested,Approved,Rejected}]
%   ...
% \end{kanban}
% \end{lstlisting}
% This option is local; it does not change the default used by the next
% environment. Letter case need not match because both declarations and task
% targets are normalized to uppercase. Other differences remain significant;
% for example, |IN REVIEW| and |REVIEW| are distinct identifiers. If a list
% repeats a name after normalization, the package reports an error, retains the
% first occurrence, and ignores the repeated occurrence. This recovery permits
% the remaining unique columns to be rendered if compilation is continued.
%
% \section{Declaring tasks}
%
% \DescribeMacro{\kanbantask}
% \SpecialUsageIndex{\kbtask}
% A task declaration has the syntax
% \begin{quote}
% |\kanbantask[|\meta{task options}|]{|\meta{column}|}{|\meta{text}|}|
% \end{quote}
% When the name is available, |\kbtask| is an exact alias with the same
% signature and behavior. The longer |\kanbantask| name is the canonical
% interface because it is less likely to conflict with another package. If
% |\kbtask| already exists when \textsf{kanban} is loaded, the package
% preserves the existing definition, issues a warning, and leaves
% |\kanbantask| fully available.
% The column must have been declared by the enclosing board. Surrounding spaces
% are removed and letters are converted to uppercase before lookup, making the
% assignment case-insensitive. The task text is ordinary LaTeX material and may
% contain emphasis, citations, and other commands suitable for a |tcolorbox|.
% Fragile or unusually complex material should be tested in the intended
% document class.
%
% The column argument may not be empty. A missing assignment or a normalized
% name that is absent from the board produces a package error. The diagnostic
% identifies the invalid input, lists the valid normalized columns, and explains
% how to correct the declaration. The invalid task is not counted or rendered,
% which leaves the remaining board data structurally consistent if compilation
% is continued after the error. A task declared outside a |kanban| environment
% similarly produces a direct error instead of modifying residual board data.
%
% For example, both tasks below are assigned to the single displayed column
% |RESEARCH|:
% \begin{lstlisting}
% \begin{kanban}[columns={Research,Review}]
%   \kanbantask{research}{Collect the evidence}
%   \kanbantask{ReSeArCh}{Document the method}
% \end{kanban}
% \end{lstlisting}
% By contrast, the following fragments are invalid:
% \begin{lstlisting}
% \kanbantask{}{A task without a target}
% \kanbantask{Archive}{A task for an undeclared stage}
% \end{lstlisting}
% For a board containing |RESEARCH| and |REVIEW|, the corresponding errors state
% that the first task has no assigned column and that |Archive| is unknown; both
% diagnostics also report |RESEARCH,REVIEW| as the valid column set.
%
% The package deliberately does not define the generic |\task| alias. That name
% is used by several document classes and packages; leaving it untouched
% prevents load-order-dependent conflicts. The namespaced |\kbtask| form
% provides a concise alternative without claiming the generic name.
%
% \subsection{Task options}
%
% \begin{description}
% \item[|prio=high|]
%   Adds a red high-priority pill, labelled \textsc{high} by default.
% \item[|prio=medium| or |prio=med|]
%   Adds an orange medium-priority pill, labelled \textsc{med} by default.
% \item[|prio=low|]
%   Adds a green low-priority pill, labelled \textsc{low} by default.
% \item[|author=|\meta{text}]
%   Places a neutral author chip at the right of the metadata row.
% \item[|deadline=|\meta{value}]
%   Adds a deadline chip beneath the metadata row. The alias |due| has exactly
%   the same meaning.
% \item[|late|]
%   Forces an existing deadline chip to use the overdue style. The option has
%   no visible effect unless |deadline| or |due| is also present.
% \item[|completed|]
%   Adds a green marker, labelled \textsc{done} by default. Completion is
%   presentational: it does not move the task or suppress its deadline.
% \end{description}
%
% Task options are local to one declaration. Unknown option names generate a
% package warning. Priority values are trimmed and compared without regard to
% letter case. An unrecognized value produces a warning and no priority pill.
%
% Several options may be combined in one declaration:
% \begin{lstlisting}
% \kanbantask[
%   prio=high,
%   author=Alice,
%   due=07/31/2026
% ]{Review}{Approve the release candidate}
% \end{lstlisting}
% The package does not infer workflow state from metadata. In particular,
% |completed| adds a marker but does not move the task to a column named
% ``Done''. The column argument remains the only source of task placement.
%
% Task text may contain ordinary inline markup:
% \begin{lstlisting}
% \kanbantask{Doing}{
%   Compare \emph{pdfLaTeX} and \emph{LuaLaTeX} output
% }
% \end{lstlisting}
% Authors should keep card text concise. Long paragraphs are valid, but they
% make narrow boards difficult to scan.
%
% \section{Deadline semantics}
%
% Deadline classification is performed when a card is rendered. By default,
% the comparison day is obtained from the TeX primitives |\year|, |\month|,
% and |\day|. Consequently, recompiling the same source on a later date may
% change a deadline chip from neutral to current or overdue. The
% |reference-date| key replaces that moving comparison day with a validated,
% fixed date when reproducible output is required.
%
% \subsection{Calendar dates}
%
% A deadline is interpreted as a calendar date when its complete expanded
% value has one of the supported numeric forms. The parser recognizes
% month-first and day-first dates with a final year, German dot notation, ISO
% notation, and slash-separated year-first notation. For example:
% \begin{lstlisting}
% \kanbantask[due=07/31/2026]{Review}{Approve the release}
% \kanbantask[due=2026-08-15]{Planned}{Archive the release}
% \end{lstlisting}
% A date before the active reference date is red, a date equal to it is orange,
% and a future date is neutral. When no reference date has been configured, the
% compilation date is the active reference. Dates are checked against the
% Gregorian month lengths, including the leap-year rule, for years from 1 to
% 9999. A numeric value that has recognized date syntax but does not denote a
% valid date produces a warning and is retained as a neutral label. Other
% descriptive values, including labels beginning with |M| or containing
% punctuation, are displayed without automatic classification; |late| may be
% used when such a label must appear overdue.
%
% \subsection{Date formats}
%
% The |date-format| key selects how every validated calendar date is printed.
% The package provides five common numeric conventions:
%
% \begin{center}
% \begin{tabular}{@{}lll@{}}
% \textbf{Key value} & \textbf{Convention} & \textbf{18 July 2026} \\
% \hline
% |us|      & month/day/year       & |07/18/2026| \\
% |english| & day/month/year       & |18/07/2026| \\
% |german|  & day.month.year       & |18.07.2026| \\
% |iso|     & year-month-day       & |2026-07-18| \\
% |ymd|     & year/month/day       & |2026/07/18| \\
% \end{tabular}
% \end{center}
%
% |us| is the initial default. Output components are zero-padded, so a valid
% input such as |2026-7-8| is not treated as ISO syntax, whereas |2026-07-08|
% is recognized and printed as |07/08/2026| under the default. Selecting a
% format follows the same three-level hierarchy as the other board settings:
% \begin{lstlisting}
% \usepackage[date-format=english]{kanban} % initial document default
% \kbset{date-format=german}               % subsequent boards
%
% \begin{kanban}[date-format=iso]           % this board only
%   ...
% \end{kanban}
% \end{lstlisting}
% The compact spelling |dateformat| is accepted as an exact key alias, although
% |date-format| is the canonical spelling used throughout this manual. Format
% values are trimmed and compared without regard to letter case. The canonical
% values in the table should be preferred in new documents. For compatibility,
% |american| maps to |us|; |uk| and |british| map to |english|; |de| and
% |deutsch| map to |german|; and |year-first| maps to |ymd|.
%
% Calendar parsing is intentionally more permissive than display formatting:
% an unambiguous date in any supported notation is accepted and normalized to
% the selected output convention. A year-last slash date can be ambiguous; for
% example, |07/08/2026| denotes either 8 July or 7 August. The |english| and
% |german| selections interpret such a value as day/month/year. The |us|,
% |iso|, and |ymd| selections use month/day/year for this otherwise ambiguous
% input. If only one ordering denotes a valid Gregorian date, that ordering is
% accepted regardless of the current selection. ISO input is consequently the
% safest form for task data exchanged between documents, because its component
% order is unambiguous. Project-month values such as |M2| and descriptive
% labels are never reformatted.
%
% The |startdate| and |reference-date| settings are normalized to an internal
% ISO representation as soon as their containing key list has been processed.
% Their values are therefore immutable snapshots: a later macro redefinition
% or board-local display-format change cannot reinterpret an inherited date.
% Within one package, setup, or board option list, |date-format| is resolved
% before either stored date is normalized, irrespective of the textual order
% of those keys. Task deadlines remain board data and are interpreted under the
% board's selected format when their cards are rendered.
%
% \subsection{Project months}
%
% A complete value of the form |M|\meta{n}, where \meta{n} is an integer from
% 1 to 999999, denotes project month \meta{n}. The upper bound keeps malformed
% input outside TeX's integer arithmetic. Values outside this range remain
% neutral labels. The initial project start date may
% be supplied as a package option:
% \begin{lstlisting}
% \usepackage[startdate=07/15/2026]{kanban}
% \end{lstlisting}
% It may be changed for subsequent boards through the common setup interface:
% \begin{lstlisting}
% \kbset{startdate=2026-07-15}
% \end{lstlisting}
% Finally, a board can override the default locally:
% \begin{lstlisting}
% \begin{kanban}[columns={Planned,Active,Done},startdate=07/15/2026]
% \end{kanban}
% \end{lstlisting}
% The precedence order is therefore board option, current |\kbset| default,
% and package option. A board-local value does not alter subsequent boards.
% An empty setup value clears the persistent start date for later boards;
% an empty board value disables an inherited start date for that board only:
% \begin{lstlisting}
% \kbset{startdate={}}
% \begin{kanban}[startdate={}]
%   ...
% \end{kanban}
% \end{lstlisting}
% Accepted start dates are stored canonically rather than as their original
% input tokens. For example, a US |07/08/2026| start date remains 8 July even
% if a later board selects |date-format=english| for its visible deadlines.
% This separation prevents a display preference from changing project-month
% boundaries.
%
% Project months begin on the day number of the start date. Thus, for a project
% beginning on 15 July, |M1| runs from 15 July until 14 August and |M2| begins
% on 15 August. If that day does not exist in a later calendar month, the
% boundary is clamped to the final valid day of that month. A project beginning
% on 31 January therefore enters |M2| on 28 February, or on 29 February in a
% leap year; its following boundary is 31 March. This rule prevents a short
% month from extending the preceding project month unintentionally.
%
% A project month earlier than the current one is overdue; the current project
% month is orange; and a later month is neutral. If no start date is available,
% the value is displayed without automatic classification. An invalid start
% date produces a package warning and does not replace the previously valid
% value.
%
% \subsection{Reproducible reference dates}
%
% The optional |reference-date| setting fixes the comparison day used by both
% calendar deadlines and project months. It follows the same package, setup,
% and board-local hierarchy as |startdate|:
% \begin{lstlisting}
% \usepackage[reference-date=2026-07-18]{kanban}
% \kbset{reference-date=2026-08-01}
%
% \begin{kanban}[reference-date=2026-08-15]
%   ...
% \end{kanban}
% \end{lstlisting}
% The key alias |today| has exactly the same meaning. An empty setup value
% restores the compilation date for later boards, while an empty board value
% restores it for that board only. Like a project start date, an accepted
% reference date is validated, expanded, and stored canonically. Invalid input
% produces a warning and preserves the preceding value. This facility is
% useful for archived reports, regression tests, and automated publication
% systems in which a rebuild should not alter status colors.
%
% \subsection{Worked deadline example}
%
% The following board uses both public key names and displays all three visual
% states with calendar dates. The manual fixes its reference date as 18 July
% 2026, so the first date lies in the past, the second equals the reference
% date, and the third lies in the future.
% \begin{lstlisting}
% \begin{kanban}[columns={Overdue,Due now,Upcoming}]
%   \kanbantask[deadline=01/01/2000]
%     {Overdue}{Resolve the review comments}
%   \kanbantask[due=07/18/2026]
%     {Due now}{Complete the current milestone}
%   \kanbantask[deadline=2099-12-31]
%     {Upcoming}{Prepare the next milestone}
% \end{kanban}
% \end{lstlisting}
% The red chip denotes a past date, the orange chip denotes the active reference
% date, and the neutral chip denotes a future date. The first and third values
% also demonstrate US and ISO calendar input. Both are printed in the selected
% format, which is US notation here because the board uses the default. The keys
% |deadline| and |due| are exact aliases; their choice does not determine the
% color. The fixed reference date keeps the rendered result stable whenever the
% manual is compiled.
%
% \begin{kanban}[columns={Overdue,Due now,Upcoming}]
%   \kanbantask[deadline=01/01/2000]
%     {Overdue}{Resolve the review comments}
%   \kanbantask[due=07/18/2026]
%     {Due now}{Complete the current milestone}
%   \kanbantask[deadline=2099-12-31]
%     {Upcoming}{Prepare the next milestone}
% \end{kanban}
%
% A complete version of this demonstration, in both light and dark mode, is
% provided as |kanban-deadlines-example.tex|.
%
% \section{Color and layout model}
%
% By default, columns draw from an eight-color palette in declaration order:
% slate, blue, teal, sage, amber, clay, rose, and plum. A selected custom
% palette may contain a different number of colors. When the number of columns
% exceeds the number of entries, the palette repeats. Priority and status
% colors are independent of the column color.
%
% In the default layout, the nominal width of each column is calculated from
% the available line width and the number of columns. Large boards can
% therefore become too narrow for useful text. In practice, three or four
% columns fit comfortably in portrait documents; wider boards often benefit
% from landscape orientation, a larger paper size, or the |break| layout.
%
% \subsection{Selecting a palette}
%
% Version 1.0.0 provides named palettes. The initial palette may be selected as
% a package option:
% \begin{lstlisting}
% \usepackage[palette=muted]{kanban}
% \end{lstlisting}
% The built-in names are |default|, |muted|, and |monochrome|. The |default|
% palette contains the original eight-color sequence.
%
% |\kbset{palette=|\meta{name}|}| changes the default for subsequent boards:
% \begin{lstlisting}
% \kbset{palette=monochrome}
% \end{lstlisting}
% A board may override that default without changing later boards:
% \begin{lstlisting}
% \begin{kanban}[columns={Todo,Doing,Done},palette=default]
%   ...
% \end{kanban}
% \end{lstlisting}
%
% \DescribeMacro{\kbdefinepalette}
% \SpecialUsageIndex{\kanbandefinepalette}
% A new palette is registered with
% \begin{quote}
% |\kanbandefinepalette{|\meta{name}|}{|\meta{color list}|}|
% \end{quote}
% where the second argument is a non-empty, comma-separated list of colors or
% \textsf{xcolor} expressions. For example:
% \begin{lstlisting}
% \definecolor{universityBlue}{HTML}{164A7B}
% \definecolor{universityGold}{HTML}{A67C00}
% \kbdefinepalette{university}{
%   universityBlue,
%   teal!70!black,
%   universityGold
% }
% \kbset{palette=university}
% \end{lstlisting}
% The unstarred form uses white heading text for every entry. If a palette
% contains both dark and light backgrounds, the starred form records an
% explicit foreground color with each background:
% \begin{quote}
% |\kanbandefinepalette*{|\meta{name}|}{|%
% \meta{background}|/|\meta{foreground}|,...}|
% \end{quote}
% For example:
% \begin{lstlisting}
% \kbdefinepalette*{accessible}{
%   facultyNavy/white,
%   facultyGold/kanbanInk,
%   facultyBerry/white
% }
% \end{lstlisting}
% The slash separates the two \textsf{xcolor} expressions. Every entry in a
% starred definition must contain exactly one slash and two non-empty values;
% otherwise the complete definition is rejected with a warning. This explicit
% pairing permits suitable heading contrast for arbitrary custom colors. It
% does not attempt to calculate or certify a contrast ratio, because
% \textsf{xcolor} expressions may depend on the surrounding document's color
% model and named-color definitions.
% When its name is available, the concise command |\kbdefinepalette| is an
% exact alias. If another package has already defined that name,
% \textsf{kanban} preserves it, reports a warning, and requires the canonical
% |\kanbandefinepalette| command. Surrounding spaces in a palette name are
% ignored, but palette names are case-sensitive. Palette registration is
% document-wide: the definition command creates or replaces a named registry
% entry, while package, setup, and board |palette| keys only select an existing
% entry.
% Defining an existing name replaces its stored color list. If that name is the
% current document default, later boards immediately use the replacement; the
% name does not need to be selected again. A board takes a local snapshot of
% the palette and its foreground list when it begins, so redefining the name
% while that board is being collected cannot alter the board already in
% progress. An unknown or empty palette produces a warning and leaves the
% current selection unchanged.
%
% A palette may contain any positive number of colors. When a board has more
% columns than the palette has entries, selection wraps cyclically. A
% three-color palette is therefore applied as colors 1, 2, 3, 1, 2, 3, and so
% forth. The built-in colors are dark enough to give white heading text a
% contrast ratio of at least 4.5:1. Authors of custom palettes remain
% responsible for preserving adequate contrast; the starred interface makes
% the foreground choice explicit for every palette entry. The following
% example defines a three-color institutional palette with dark text on its
% lighter gold background. Since the board has four columns, its
% fourth heading reuses the first palette entry.
% \begin{lstlisting}
% \definecolor{facultyNavy}{HTML}{244A6C}
% \definecolor{facultyGold}{HTML}{C69A32}
% \definecolor{facultyBerry}{HTML}{843D60}
% \kbdefinepalette*{faculty}{
%   facultyNavy/white,
%   facultyGold/kanbanInk,
%   facultyBerry/white
% }
%
% \begin{kanban}[
%   columns={Proposed,Selected,In progress,Published},
%   palette=faculty
% ]
%   \kanbantask[deadline=12/31/2099]
%     {Proposed}{Collect research questions}
%   \kanbantask[prio=high,deadline=01/01/2000]
%     {Selected}{Approve the study design}
%   \kanbantask[author=Alice,
%     due=07/18/2026]
%     {In progress}{Complete the analysis}
%   \kanbantask[completed]{Published}{Deposit the article}
% \end{kanban}
% \end{lstlisting}
% It produces the following board:
%
% \definecolor{kbManualNavy}{HTML}{244A6C}
% \definecolor{kbManualGold}{HTML}{C69A32}
% \definecolor{kbManualBerry}{HTML}{843D60}
% \kbdefinepalette*{manualfaculty}{
%   kbManualNavy/white,
%   kbManualGold/kanbanInk,
%   kbManualBerry/white
% }
% \begin{kanban}[
%   columns={Proposed,Selected,In progress,Published},
%   palette=manualfaculty
% ]
%   \kanbantask[deadline=12/31/2099]
%     {Proposed}{Collect research questions}
%   \kanbantask[prio=high,deadline=01/01/2000]
%     {Selected}{Approve the study design}
%   \kanbantask[author=Alice,
%     due=07/18/2026]
%     {In progress}{Complete the analysis}
%   \kanbantask[completed]{Published}{Deposit the article}
% \end{kanban}
%
% The complete compilable source is |kanban-palette-example.tex|.

% \subsection{Column geometry}
%
% The |columnsep| key accepts either a non-negative explicit TeX dimension,
% such as |2mm| or |0.5em|, or a single registered dimension such as
% |\linewidth|. It controls the distance between adjacent columns, and its
% default value is |2mm|. Malformed values are rejected with a package warning
% before TeX's dimension scanner is entered, so they do not produce secondary
% ``missing number'' or ``illegal unit'' errors. An initial value may be set as
% a package option:
% \begin{lstlisting}
% \usepackage[columnsep=3mm]{kanban}
% \end{lstlisting}
% It may be changed for subsequent boards with the shared setup command:
% \begin{lstlisting}
% \kbset{columnsep=1.5mm}
% \end{lstlisting}
% A single board may override the current default without affecting later
% boards:
% \begin{lstlisting}
% \begin{kanban}[columns={Todo,Doing,Done},columnsep=4mm]
%   ...
% \end{kanban}
% \end{lstlisting}
% The same value is applied to the fixed layout and to the column separation
% used by the |break| layout. In the fixed layout the package subtracts the
% total separation from the available line width before dividing the remaining
% width among the columns. Each task card then occupies the complete available
% width of its column body.
%
% The value must also be feasible for the resolved number of columns. For
% \meta{n} columns, the product of \meta{n} minus one and |columnsep| must be
% smaller than |\linewidth|. If the total gap is equal to or wider than the
% available line width, the package reports an error and uses |0pt| for that
% board. This recovery value preserves positive column width while leaving the
% document-wide default unchanged.
%
% The fixed layout is constructed as a two-row raster. All heading cells use
% the height of the tallest heading, including headings that wrap onto more
% than one line. All column bodies use the height of the body containing the
% most material, so empty or short columns extend to the same bottom edge. The
% column body itself is transparent in both visual modes; only its bottom rule
% remains visible, while the left and right rules are omitted. Task cards keep
% their light or dark background, border, and colored accent.
%
% Equal-height measurements are stored in the document's auxiliary file.
% Authors should therefore compile the document twice after changing a column
% heading or task list. In the page-breakable layout, heading heights are also
% synchronized at the beginning of the board, but column bodies remain
% independent so that they can continue naturally across pages. The initial
% headings are not repeated on continuation pages.
%
% The following example uses concise headings and includes an empty column. The
% shorter labels are preferable in narrow boards because they keep the workflow
% stages visually compact; longer headings remain supported and are equalized
% automatically when wrapping is unavoidable.
% \begin{lstlisting}
% \begin{kanban}[columns={Queued,Review,Complete}]
%   \kanbantask[prio=high]{Review}
%     {Prepare the response}
%   \kanbantask{Review}
%     {Check the revised figures}
%   \kanbantask[completed]{Complete}{Approve the revision}
% \end{kanban}
% \end{lstlisting}
% \Needspace{10\baselineskip}
% It produces a board whose headings and bottom edges remain aligned:
%
% \begin{kanban}[columns={Queued,Review,Complete}]
%   \kanbantask[prio=high]{Review}
%     {Prepare the response}
%   \kanbantask{Review}
%     {Check the revised figures}
%   \kanbantask[completed]{Complete}{Approve the revision}
% \end{kanban}
%
% \subsection{Light and dark modes}
%
% The default |light| mode uses dark text on white cards, neutral gray metadata
% chips, and white headings on the selected palette colors. The |dark| mode
% uses light text on dark cards. Author chips and neutral |deadline| or |due|
% chips deliberately retain their light-mode colors in both modes, giving
% metadata a stable visual identity. Automatic status colors remain semantic:
% overdue values are red and values due now are amber in either mode. In both
% modes the area behind the cards is transparent, so the surrounding page
% remains visible in the vertical space between cards and at the foot of a
% synchronized column. Palette colors continue to identify columns in both
% modes.
%
% The initial mode may be selected as a package option:
% \begin{lstlisting}
% \usepackage[mode=dark]{kanban}
% \end{lstlisting}
% It can be changed for subsequent boards through the common setup interface:
% \begin{lstlisting}
% \kbset{mode=dark}
% \end{lstlisting}
% A board may then override that default locally:
% \begin{lstlisting}
% \begin{kanban}[columns={Todo,Doing,Done},mode=light]
%   ...
% \end{kanban}
% \end{lstlisting}
%
% Mode selection affects only the board. It deliberately does not call
% |\pagecolor| and does not alter the surrounding document text. A dark board
% can therefore be embedded safely in an otherwise light document. Unknown
% mode names produce a warning and leave the current selection unchanged.
%
% The examples section below compares the two modes with identical column
% geometry and explicit deadline states. A separate custom-palette example is
% provided as |kanban-palette-example.tex|.
%
% \subsection{Labels and public style hooks}
%
% The visible semantic labels can be adapted without redefining internal
% commands. The |labels| key accepts the nested keys |high|, |medium|, |low|,
% and |done|. It follows the same package, setup, and board-local precedence as
% the other configuration keys:
% \begin{lstlisting}
% \usepackage[
%   labels={high=URGENT,medium=NORMAL,low=ROUTINE,done=FINISHED}
% ]{kanban}
%
% \kbset{labels={done=APPROVED}}
%
% \begin{kanban}[labels={high=IMMEDIATE}]
%   ...
% \end{kanban}
% \end{lstlisting}
% A board takes a snapshot of all four labels when it begins. An omitted label
% keeps its current value; an empty value is rejected with a warning. Label
% text is ordinary short LaTeX material, although concise uppercase terms are
% normally easiest to read inside the small chips.
%
% Seven namespaced \textsf{tcolorbox} styles provide a controlled extension
% point for typographic and spacing changes:
% \begin{quote}
% |kanban/header|, |kanban/card|, |kanban/column-body|;\\
% |kanban/priority-chip|, |kanban/author-chip|;\\
% |kanban/deadline-chip|, |kanban/completed-chip|.
% \end{quote}
% The styles are initially empty and are applied after the structural package
% defaults. For example:
% \begin{lstlisting}
% \tcbset{
%   kanban/header/.append style={fontupper=\bfseries\sffamily},
%   kanban/card/.append style={arc=1mm,outer arc=1mm},
%   kanban/author-chip/.append style={left=1.5mm,right=1.5mm}
% }
% \end{lstlisting}
% These hooks are document-wide \textsf{tcolorbox} settings. The header, card,
% and chip styles apply to both rendering strategies. |kanban/column-body|
% applies only to the fixed renderer because the page-breakable renderer leaves
% each \textsf{paracol} stream without an enclosing body box. Authors should
% preserve |width=\linewidth| for cards and should test changes to padding in
% narrow columns. The named colors |kanbanHigh|, |kanbanMed|, and |kanbanLow|
% remain available to advanced users, but changing either foreground or
% background colors also transfers responsibility for readable contrast to the
% document author. In the defaults, the orange medium and due-now chips use
% dark ink; the red and green semantic chips use white. Together with the
% revised built-in palettes, these pairs meet a 4.5:1 reference contrast ratio
% for ordinary text.
%
% \section{Diagnostics and limitations}
%
% Unknown setup, board, and task options generate package warnings. Unknown
% palettes and modes preserve the current selection; an unknown priority is
% omitted. Empty semantic labels preserve the preceding value. A numeric
% deadline that resembles a calendar date or project month but is outside the
% validated range is displayed neutrally and produces a warning; descriptive
% labels remain neutral without a diagnostic. A malformed or negative
% |columnsep| value is rejected with one package warning and preserves the
% current value; malformed input is intercepted before TeX emits a low-level
% number or unit error. A non-negative separation whose combined gaps leave no column width
% produces an error and is replaced by |0pt| for the affected board. A task with an
% empty column argument, a task outside a board, a task assigned to an
% undeclared normalized name, or a column list containing a duplicate after
% case normalization generates a package error with corrective help. Duplicate
% names recover by keeping the first normalized occurrence and discarding later
% occurrences. Nested
% |kanban| environments are also rejected because each board owns one active
% collection context. The nested environment produces one error; task
% declarations inside it are then ignored without a cascade of secondary
% ``outside a board'' errors. Invalid tasks are neither counted nor rendered. This
% distinction reserves errors for
% declarations that cannot be interpreted unambiguously while allowing harmless
% option diagnostics to remain recoverable warnings.
% The optional short commands |\kbtask|, |\kbset|, and |\kbdefinepalette| are
% installed only when their names are unused. A pre-existing definition is
% never overwritten; the canonical long command remains available in every
% case.
%
% Authors should also observe the following constraints.
% \begin{itemize}
% \item The resolved default or explicit list must contain at least one column.
%   The built-in default satisfies this requirement, and an empty value passed
%   to the |columns| setup key is rejected.
% \item Column identifiers are converted to uppercase and compared without
%   regard to letter case. They must not contain commas and must remain unique
%   after normalization.
% \item Mode and priority values are trimmed and compared without regard to
%   letter case. Date-format values and their documented aliases follow the
%   same normalization rule. Palette names remain case-sensitive.
% \item Column separation must be a non-negative explicit dimension or a
%   single registered dimension, and the sum of all inter-column gaps must be
%   smaller than the available line width.
% \item Calendar dates, project start dates, and reference dates may use US or English
%   slash notation with a final year, German dot notation, ISO hyphen
%   notation, or slash-separated year-first notation. Ambiguous year-last
%   slash dates follow the selected date format. Calendar dates are validated
%   for years from 1 to 9999 and are normalized for display.
% \item Only complete calendar dates and values |M1| through |M999999| are
%   classified automatically. All other deadline values are treated as labels
%   and remain neutral unless |late| is present. The |M| in a project-month
%   value is uppercase and case-sensitive.
% \item The fixed raster layout cannot split an individual board across pages;
%   use |break| for long boards.
% \item The |break| layout prints synchronized headings once at the beginning
%   of the board. It does not repeat them on continuation pages because the
%   supported \textsf{paracol} interface provides independent column streams,
%   not a repeating per-column heading row.
% \item The |break| layout should not be used inside a
%   \textsf{pdflscape} |landscape| environment. Use the fixed layout for a
%   landscape board and ensure that the complete board fits on that page.
% \item Task counts include every successfully assigned task, irrespective of
%   its completion status.
% \end{itemize}
%
% \section{Further usage examples}
%
% The distribution contains eight independently compilable example documents.
% The general overview is |kanban-example.tex|; the remaining files isolate
% deadline states, date formats, custom palettes, visual modes, column
% separation, page breaking, and landscape composition.
% This section reproduces the essential source for each case and, where visual
% behaviour is central, places the generated output directly below it. The
% resulting comparison makes the effect of each option observable without
% requiring the reader to compile a separate file.
%
% \subsection{Reusable document setup}
%
% A document containing several related boards should establish common values
% at the highest appropriate configuration level. Package options provide the
% initial policy, while environment options record deliberate exceptions. The
% following example keeps the column list, palette, mode, project start date,
% reference date, and column separation consistent across the document, but gives the second board
% a local column list and a different mode:
% \begin{lstlisting}
% \documentclass{article}
% \usepackage[
%   columns={Planned,Active,Complete},
%   palette=muted,
%   mode=light,
%   startdate=07/01/2026,
%   reference-date=2026-07-18,
%   columnsep=2mm
% ]{kanban}
%
% \begin{document}
%
% \section{Research tasks}
% \begin{kanban}
%   \kanbantask[prio=high,author=Alice,due=M1]
%     {Planned}{Validate the data set}
%   \kanbantask[author=Bob,due=M2]
%     {Active}{Draft the results section}
%   \kanbantask[completed]
%     {Complete}{Approve the analysis plan}
% \end{kanban}
%
% \section{Archival tasks}
% \begin{kanban}[columns={Queued,Processing,Stored},mode=dark]
%   \kanbantask{Queued}{Deposit the source data}
%   \kanbantask[author=Carol]{Processing}{Create metadata records}
% \end{kanban}
%
% \end{document}
% \end{lstlisting}
% The first board inherits every package default, including the configured
% column list. The second board inherits the |muted| palette, project start
% date, and column separation, while replacing its columns and visual mode for
% that board only. This pattern makes document-wide conventions explicit and
% keeps local deviations easy to identify during review.
%
% \Needspace{18\baselineskip}
% \subsection{Selecting a date format}
%
% A shared date convention should normally be established at package or setup
% level. A board-local selection is useful when a report reproduces material
% prepared for a different audience. The following source uses ISO input for
% every task because that form is unambiguous, then requests each available
% output convention independently:
% \begin{lstlisting}
% \begin{kanban}[columns={Upcoming},date-format=us]
%   \kanbantask[due=2099-07-18]
%     {Upcoming}{Prepare the scheduled release}
% \end{kanban}
%
% \begin{kanban}[columns={Upcoming},date-format=english]
%   \kanbantask[due=2099-07-18]
%     {Upcoming}{Prepare the scheduled release}
% \end{kanban}
%
% % The same structure also accepts german, iso, and ymd.
% \end{lstlisting}
% Only the printed date changes; validation and overdue classification use the
% same Gregorian components in all five cases.
%
% \paragraph{Rendered output.}
% The boards below render the same future date. The labels state the component
% order explicitly so that the result does not depend on national terminology.
%
% \noindent
% \kanbanManualDateFormatBoard{us}{US: MM/DD/YYYY}
% \hfill
% \kanbanManualDateFormatBoard{english}{English: DD/MM/YYYY}
%
% \medskip
%
% \noindent
% \kanbanManualDateFormatBoard{german}{German: DD.MM.YYYY}
% \hfill
% \kanbanManualDateFormatBoard{iso}{ISO: YYYY-MM-DD}
%
% \medskip
%
% \noindent
% \kanbanManualDateFormatBoard{ymd}{Year first: YYYY/MM/DD}
%
% The complete compilable source is |kanban-date-formats-example.tex|. It also
% shows how the five compact boards can be arranged in minipages without
% changing their local configuration scope.
%
% \Needspace{18\baselineskip}
% \subsection{Light and dark boards in one document}
%
% Visual mode is a property of the cards rather than of the page. Selecting
% |dark| therefore changes card backgrounds, text, and borders
% without changing the surrounding document or the transparent column body.
% The source below keeps the column structure constant so that the change in
% visual treatment can be assessed directly. Author chips and neutral deadline
% chips keep the same light treatment in both boards; red and amber deadline
% states retain their semantic colors.
%
% \begin{samepage}
% \begin{lstlisting}
% \begin{kanban}[columns={Planned,Active,Complete},mode=light]
%   \kanbantask[prio=high,author=Alice,deadline=01/01/2000]
%     {Planned}{Confirm the requirements}
%   \kanbantask[author=Bob,
%     due=07/18/2026]
%     {Active}{Implement the package interface}
%   \kanbantask[completed,deadline=12/31/2099]
%     {Complete}{Create the test files}
% \end{kanban}
%
% \bigskip
%
% \begin{kanban}[columns={Planned,Active,Complete},mode=dark]
%   \kanbantask[prio=medium,author=Carol,deadline=01/01/2000]
%     {Planned}{Prepare the release notes}
%   \kanbantask[author=David,
%     due=07/18/2026]
%     {Active}{Review the documentation}
%   \kanbantask[completed,deadline=12/31/2099]
%     {Complete}{Select the license}
% \end{kanban}
% \end{lstlisting}
% \end{samepage}
% A document-wide mode may instead be selected as a package option or with
% |\kbset|. The environment option remains the appropriate form when the change
% applies to one board only.
%
% \Needspace{14\baselineskip}
% \paragraph{Rendered output.}
% The following pair is generated from the preceding source. In particular,
% the dark board retains the selected column palette while using transparent
% column bodies, dark cards, light foreground text, and light neutral author
% and future-deadline chips matching those in the light board. Red and amber
% deadline chips continue to communicate status independently of visual mode.
%
% \begin{kanban}[columns={Planned,Active,Complete},mode=light]
%   \kanbantask[prio=high,author=Alice,deadline=01/01/2000]
%     {Planned}{Confirm the requirements}
%   \kanbantask[author=Bob,
%     due=07/18/2026]
%     {Active}{Implement the package interface}
%   \kanbantask[completed,deadline=12/31/2099]
%     {Complete}{Create the test files}
% \end{kanban}
%
% \medskip
%
% \begin{kanban}[columns={Planned,Active,Complete},mode=dark]
%   \kanbantask[prio=medium,author=Carol,deadline=01/01/2000]
%     {Planned}{Prepare the release notes}
%   \kanbantask[author=David,
%     due=07/18/2026]
%     {Active}{Review the documentation}
%   \kanbantask[completed,deadline=12/31/2099]
%     {Complete}{Select the license}
% \end{kanban}
%
% The complete compilable source is provided as |kanban-modes-example.tex|.
%
% \subsection{Changing the gap between columns}
%
% Column separation follows the same three-level configuration hierarchy as
% the palette and mode. A package option establishes the initial value,
% |\kbset| changes the value for subsequent boards, and an environment option
% applies to one board. The following fragment demonstrates the scope of each
% form:
%
% \begin{samepage}
% \begin{lstlisting}
% \usepackage[columnsep=3mm]{kanban}
%
% % Uses the 3 mm package default.
% \begin{kanban}[columns={Todo,Doing,Done}]
%   \kanbantask[deadline=01/01/2000]{Todo}{Define the interface}
%   \kanbantask[due=07/18/2026]
%     {Doing}{Write the implementation}
%   \kanbantask[deadline=12/31/2099,completed]{Done}{Run the tests}
% \end{kanban}
%
% % Changes the default for subsequent boards.
% \kbset{columnsep=1mm}
% \begin{kanban}[columns={Todo,Doing,Done}]
%   \kanbantask[deadline=01/01/2000]{Todo}{Define the interface}
%   \kanbantask[due=07/18/2026]
%     {Doing}{Write the implementation}
%   \kanbantask[deadline=12/31/2099,completed]{Done}{Run the tests}
% \end{kanban}
%
% % Uses 4 mm here; the document default remains 1 mm.
% \begin{kanban}[columns={Todo,Doing,Done},columnsep=4mm]
%   \kanbantask[deadline=01/01/2000]{Todo}{Define the interface}
%   \kanbantask[due=07/18/2026]
%     {Doing}{Write the implementation}
%   \kanbantask[deadline=12/31/2099,completed]{Done}{Run the tests}
% \end{kanban}
% \end{lstlisting}
% \end{samepage}
% The value must be a nonnegative explicit dimension, such as |2mm|, or one
% registered dimension, such as |\linewidth|. Increasing it strengthens the
% visual separation between workflow stages, but reduces the width available
% to every task card.
%
% \paragraph{Rendered output.}
% In the comparison below, palette, headings, and task content are held
% constant. Only |columnsep| changes. Board-local values are used here so that
% the manual's configuration remains unchanged after the demonstration.
%
% \noindent\textit{Column separation: 3 mm.}
% \begin{kanban}[columns={Todo,Doing,Done},columnsep=3mm]
%   \kanbantask[prio=high,deadline=01/01/2000]
%     {Todo}{Define the interface}
%   \kanbantask[author=Alice,
%     due=07/18/2026]
%     {Doing}{Write the implementation}
%   \kanbantask[completed,deadline=12/31/2099]{Done}{Run the tests}
% \end{kanban}
%
% \medskip
% \noindent\textit{Column separation: 1 mm.}
% \begin{kanban}[columns={Todo,Doing,Done},columnsep=1mm]
%   \kanbantask[prio=high,deadline=01/01/2000]
%     {Todo}{Define the interface}
%   \kanbantask[author=Alice,
%     due=07/18/2026]
%     {Doing}{Write the implementation}
%   \kanbantask[completed,deadline=12/31/2099]{Done}{Run the tests}
% \end{kanban}
%
% \medskip
% \Needspace{8\baselineskip}
% \noindent\textit{Column separation: 4 mm.}
% \begin{kanban}[columns={Todo,Doing,Done},columnsep=4mm]
%   \kanbantask[prio=high,deadline=01/01/2000]
%     {Todo}{Define the interface}
%   \kanbantask[author=Alice,
%     due=07/18/2026]
%     {Doing}{Write the implementation}
%   \kanbantask[completed,deadline=12/31/2099]{Done}{Run the tests}
% \end{kanban}
%
% The complete source is |kanban-columnsep-example.tex|.
%
% \Needspace{10\baselineskip}
% \subsection{A page-breakable board}
%
% Long task lists should use the |break| option. In this layout each Kanban
% column is managed as a \textsf{paracol} column and may continue onto a later
% page. The fixed layout, by contrast, treats the complete board as one
% unbreakable object. The compact fragment below shows the public interface;
% the rendered demonstration that follows deliberately contains more cards so
% that a real page transition occurs. The synchronized headings appear once at
% the beginning of the board; continuation pages contain only the remaining
% cards. This behavior follows the public \textsf{paracol} column-stream model
% and avoids dependencies on its private output routine.
%
% \begin{samepage}
% \begin{lstlisting}
% \kbset{palette=default,mode=light,columnsep=2mm}
% \begin{kanban}[
%   columns={Backlog,In progress,Review,Done},break
% ]
%   \kanbantask[prio=high,author=Alice,deadline=01/01/2000]
%     {Backlog}{Finalize the interface}
%   \kanbantask[prio=medium,author=Bob,
%     due=07/18/2026]
%     {In progress}{Write regression tests}
%   \kanbantask[author=Carol,deadline=12/31/2099]
%     {Review}{Proofread the manual}
%   \kanbantask[completed,due=12/31/2099]
%     {Done}{Select the package license}
% \end{kanban}
% \end{lstlisting}
% \end{samepage}
% Task declarations need not follow the order of the columns. The package
% collects each task under its named column and preserves source order within
% that column. Since column bodies paginate independently, their bottom edges
% are not synchronized in this layout.
%
% \paragraph{Rendered output.}
% The next board begins on a fresh page and contains enough material to continue
% onto the following page. Its headings are printed once at the beginning;
% subsequent cards remain in their respective column streams.
%
% \clearpage
% \begin{kanban}[
%   columns={Backlog,In progress,Review},break,columnsep=2mm
% ]
%   \kanbanManualBreakDeadlineCard{Backlog}{Alice}
%     {deadline=01/01/2000}{Confirm the public interface}
%   \kanbanManualBreakCard{Backlog}{Bob}{Collect compatibility requirements}
%   \kanbanManualBreakCard{Backlog}{Carol}{Review package dependencies}
%   \kanbanManualBreakCard{Backlog}{David}{Define the project date policy}
%   \kanbanManualBreakCard{Backlog}{Alice}{Select the release palette}
%   \kanbanManualBreakCard{Backlog}{Bob}{Identify supported engines}
%   \kanbanManualBreakCard{Backlog}{Carol}{Plan the documentation structure}
%   \kanbanManualBreakCard{Backlog}{David}{Prepare the release checklist}
%   \kanbanManualBreakCard{Backlog}{Alice}{Confirm the archive structure}
%   \kanbanManualBreakCard{Backlog}{Bob}{Review installation instructions}
%   \kanbanManualBreakCard{Backlog}{Carol}{Assess documented limitations}
%   \kanbanManualBreakCard{Backlog}{David}{Schedule the release review}
%
%   \kanbanManualBreakDeadlineCard{In progress}{Alice}
%     {due=07/18/2026}
%     {Implement configuration keys}
%   \kanbanManualBreakCard{In progress}{Bob}{Apply the visual modes}
%   \kanbanManualBreakCard{In progress}{Carol}{Synchronize header heights}
%   \kanbanManualBreakCard{In progress}{David}{Equalize fixed column bodies}
%   \kanbanManualBreakCard{In progress}{Alice}{Remove outer side rules}
%   \kanbanManualBreakCard{In progress}{Bob}{Extend the user manual}
%   \kanbanManualBreakCard{In progress}{Carol}{Compile example documents}
%   \kanbanManualBreakCard{In progress}{David}{Add regression assertions}
%   \kanbanManualBreakCard{In progress}{Alice}{Run the pdfTeX checks}
%   \kanbanManualBreakCard{In progress}{Bob}{Run the LuaTeX checks}
%   \kanbanManualBreakCard{In progress}{Carol}{Generate the release archive}
%   \kanbanManualBreakCard{In progress}{David}{Inspect the rendered output}
%
%   \kanbanManualBreakDeadlineCard{Review}{Alice}
%     {deadline=12/31/2099}{Inspect the generated style file}
%   \kanbanManualBreakCard{Review}{Bob}{Compare regression logs}
%   \kanbanManualBreakCard{Review}{Carol}{Proofread the public interface}
%   \kanbanManualBreakCard{Review}{David}{Examine color contrast}
%   \kanbanManualBreakCard{Review}{Alice}{Validate the documented source}
%   \kanbanManualBreakCard{Review}{Bob}{Review the CTAN metadata}
%   \kanbanManualBreakCard{Review}{Carol}{Check the license statements}
%   \kanbanManualBreakCard{Review}{David}{Inspect the manual PDF}
%   \kanbanManualBreakCard{Review}{Alice}{Test the installation procedure}
%   \kanbanManualBreakCard{Review}{Bob}{Verify all example files}
%   \kanbanManualBreakCard{Review}{Carol}{Confirm the change history}
%   \kanbanManualBreakCard{Review}{David}{Approve the release candidate}
% \end{kanban}
%
% The continuation demonstrates why body heights remain independent when
% |break| is active. The complete standalone example is
% |kanban-break-example.tex|.
%
% \subsection{A fixed board on a landscape page}
%
% A landscape page increases |\linewidth| and is therefore useful when a
% workflow has five or more stages that would be too narrow in portrait
% orientation. The \textsf{pdflscape} package establishes the rotated page and
% records the appropriate PDF rotation metadata; the board itself requires no
% rotation commands. Page margins should be chosen before the board is opened,
% and |columnsep| can then be reduced modestly if the available width is still
% insufficient.
%
% The board must use the fixed layout. The |break| option is deliberately
% absent because its \textsf{paracol}-based pagination is intended for ordinary
% portrait text pages and is not compatible with this use. A fixed board is
% kept as one indivisible object, including its header and body rows. It will
% therefore move to the next page when the remaining vertical space is too
% small. If the complete board is taller than one landscape page, the content
% should be reduced or presented as a separate page-breakable portrait board.
%
% Concise stage names are especially important when six columns share the page.
% The example uses |Build| and |Docs| as familiar abbreviations for
% implementation and documentation. Less familiar abbreviations should be
% defined in the surrounding text so that the board remains understandable
% without relying on source-code comments.
%
% \begin{samepage}
% \begin{lstlisting}
% \documentclass[a4paper]{article}
% \usepackage[margin=18mm]{geometry}
% \usepackage{pdflscape}
% \usepackage[columnsep=2mm]{kanban}
%
% \begin{document}
% \begin{landscape}
% \section*{Release workflow}
%
% % Deliberately omit the break option in a landscape environment.
% \begin{kanban}[
%   columns={Ideas,Selected,Build,Review,Docs,Released},
%   mode=light
% ]
%   \kanbantask[deadline=01/01/2000]
%     {Ideas}{Collect feature requests}
%   \kanbantask{Ideas}{Record user feedback}
%   \kanbantask[prio=high,
%     due=07/18/2026]
%     {Selected}{Approve the release scope}
%   \kanbantask[author=Alice,deadline=12/31/2099]
%     {Build}{Complete the changes}
%   \kanbantask[author=Bob]{Review}{Run the acceptance review}
%   \kanbantask[author=Carol]{Docs}{Update the user manual}
%   \kanbantask[completed]{Released}{Publish the archive}
% \end{kanban}
% \end{landscape}
% \end{document}
% \end{lstlisting}
% \end{samepage}
% The supporting package is needed only by a document that requests rotated
% pages; it is not a dependency of \textsf{kanban}. The excerpt shows the
% structural pattern. The standalone example adds three tasks to every stage so
% that the use of the available landscape area can be assessed realistically.
%
% \paragraph{Rendered output.}
% The next page is generated directly in a |landscape| environment. The board
% uses six fixed columns, short one-line headings, three cards per column, and
% no |break| option. Equal body heights make the workflow read as one coherent
% table even though individual task cards have different amounts of metadata.
%
% \begin{landscape}
% \section*{Landscape release workflow}
% \begin{kanban}[
%   columns={Ideas,Selected,Build,Review,Docs,Released},
%   mode=light,
%   columnsep=2mm
% ]
%   \kanbantask[author=Alice,deadline=01/01/2000]
%     {Ideas}{Collect feature requests}
%   \kanbantask[author=Bob]
%     {Ideas}{Record user feedback}
%   \kanbantask[author=Carol]
%     {Ideas}{Identify compatibility risks}
%
%   \kanbantask[prio=high,author=Bob,
%     due=07/18/2026]
%     {Selected}{Approve the release scope}
%   \kanbantask[author=Carol]
%     {Selected}{Confirm supported engines}
%   \kanbantask[author=David]
%     {Selected}{Assign milestone owners}
%
%   \kanbantask[author=Carol,deadline=12/31/2099]
%     {Build}{Complete the package changes}
%   \kanbantask[author=David]
%     {Build}{Update date validation}
%   \kanbantask[author=Alice]
%     {Build}{Generate the style file}
%
%   \kanbantask[author=David]
%     {Review}{Run the technical review}
%   \kanbantask[author=Alice]
%     {Review}{Compare engine output}
%   \kanbantask[author=Bob]
%     {Review}{Check regression logs}
%
%   \kanbantask[author=Alice]
%     {Docs}{Update the user manual}
%   \kanbantask[author=Bob]
%     {Docs}{Compile usage examples}
%   \kanbantask[author=Carol]
%     {Docs}{Verify CTAN metadata}
%
%   \kanbantask[completed]
%     {Released}{Publish the source archive}
%   \kanbantask[author=Alice,completed]
%     {Released}{Tag version 1.0.0}
%   \kanbantask[author=Bob,completed]
%     {Released}{Announce the release}
% \end{kanban}
% \end{landscape}
%
% The output confirms that the fixed layout uses the landscape line width while
% retaining synchronized headings and column bodies. The task count in each
% heading is derived from the declarations, and the palette continues cyclically
% if a later board contains more stages than palette entries. A complete version
% is provided as |kanban-landscape-example.tex|.
%
% \section{License and maintenance}
%
% The package is distributed under the LaTeX Project Public License, version
% 1.3c or, at the user's option, any later version. Its LPPL status is
% ``maintained'', and the Current Maintainer is Andreas Willibald Weber. Defect reports
% should include a minimal complete document, the relevant log messages, and
% the versions of LaTeX and \textsf{kanban} used to reproduce the problem. The
% package information page is \url{https://ctan.org/pkg/kanban}. Development
% takes place at
% \url{https://www.github.com/andreaswillibaldweber/ctan-kanban}; defects and
% focused enhancement proposals should be submitted through the public issue
% tracker at
% \url{https://www.github.com/andreaswillibaldweber/ctan-kanban/issues}.
% The source archive includes |CONTRIBUTING.md| for contributors and
% |RELEASING.md| for the maintainer's reproducible publication procedure.
%
% \StopEventually{\clearpage\PrintChanges\PrintIndex}
%
% \section{Implementation}
%
% \subsection{Public command contracts}
%
% The implementation distinguishes persistent defaults, one-board state, and
% task data. The |kanban| environment establishes a collection scope, copies
% the persistent defaults, applies its local options, and selects one of the
% two renderers. It is the only public operation that creates board state.
% |\kanbantask|, or its exact short alias |\kbtask| when available, may mutate
% that state only while a board is active; otherwise it reports an error and
% discards the declaration. A nested board is rejected at
% its entry point, and its task declarations enter a temporary discard state so
% that the primary structural error is not obscured by secondary diagnostics.
%
% |\kanbansetup| changes persistent defaults for boards that begin later. When
% the short name is unused, |\kbset| is defined by control-sequence equivalence,
% not by a forwarding wrapper. The same rule applies to
% |\kanbandefinepalette| and |\kbdefinepalette|, and to |\kanbantask| and
% |\kbtask|. Existing short definitions are preserved and diagnosed. Palette
% definition changes the registry; selection
% changes the current palette name, its background list, and any paired
% heading-foreground list. Redefining the selected name refreshes the
% persistent lists for later boards, whereas a board already in progress keeps
% its local snapshot. No generic |\task| command is defined,
% because that generic name has a material risk of collision with classes and
% unrelated packages.
%
% \subsection{Configuration and data collection}
%
% A board begins by copying the current document defaults for columns, palette,
% visual mode, date format, project start date, reference date, semantic labels,
% and column separation into board-local state.
% Its optional key list is then applied to that state. In particular, the
% |columns| key replaces the copied column list for one board without changing
% later boards. Every column name is trimmed and converted to uppercase. A
% board-local LaTeX3 property list maps each normalized name to a consecutive
% numeric identifier. Task lists and counters use the numeric identifier and a
% monotonically increasing board identifier, so user-provided names never form
% part of an internal control-sequence name. Duplicate normalized names are
% diagnosed before rendering. Column lists, palette backgrounds, and optional
% paired foregrounds are stored as LaTeX3 comma lists. Palette entries are
% selected cyclically by column position, so a palette may contain either
% fewer or more colors than a particular board requires. Configuration is
% parsed by the namespaced LaTeX3 key families |kanban/task|, |kanban/board|,
% and |kanban/setup|, with subordinate label families. Keeping these roles
% explicit avoids leaking a generic key family into the document namespace.
% Start and reference dates are deferred until the complete surrounding key
% list has resolved its date format, then validated, expanded, and stored in
% canonical ISO form. Package options are accumulated and passed through the
% same setup transaction rather than maintaining a second parser.
%
% Task declarations do not typeset cards immediately. Instead, each declaration
% normalizes and validates the target column, increments its counter, and
% appends a protected call to the card renderer to that column's list. Empty and
% unknown targets produce explicit errors and are not appended. Task text and
% metadata remain unexpanded while they are stored. Collection is global only
% for the lifetime of the active board, which means a declaration inside an
% ordinary TeX group is retained. The board identifier isolates that temporary
% state, and the generated names, colors, heading foregrounds, lists, and
% counters are explicitly undefined after rendering. The lists are expanded
% only when the environment closes. This
% deferred model preserves declaration order within each column and allows
% tasks for different columns to be interleaved freely in the document source.
% After the columns are resolved, the package also evaluates the complete gap
% width. Dimension syntax is checked before assignment, preventing malformed
% input from reaching TeX's number scanner. A value that would leave no positive
% width for cards is diagnosed once
% and replaced locally by |0pt| before either renderer is entered.
%
% \subsection{Rendering strategies}
%
% The fixed renderer creates two raster rows: one containing every heading and
% one containing every transparent column body. Row-wise equalization gives all
% headings the height of the tallest heading and all bodies the height of the
% fullest body. The complete raster is enclosed in a minipage. This makes the
% header and body rows indivisible, preventing an orphaned header row at the
% foot of a page. A fixed board that does not fit in the remaining space moves
% to the following page; a board taller than the page must instead use the
% page-breakable layout.
%
% When |break| is selected, the package constructs parallel \textsf{paracol}
% columns. Header heights are synchronized through a board-specific equal-height
% group, while the card lists remain outside an enclosing fixed-height box. Each
% column may therefore continue naturally on a later page. Headers are emitted
% once before the streams diverge and are not repeated on continuation pages;
% this uses only the supported \textsf{paracol} interface. The two renderers
% share the same card command and visual styles, but deliberately use different
% body-height and pagination rules.
%
% \subsection{Metadata and validation}
%
% Calendar and project-month classification is performed when a card is
% rendered. The deadline value is expanded into a temporary token list and is
% matched against exact numeric forms. All five supported calendar notations
% are reduced to the same year, month, and day components. Ambiguous year-last
% slash values are interpreted according to the active date-format snapshot,
% with a validity-based swap when only the opposite ordering forms a real
% Gregorian date. The components are checked against month lengths and leap-year
% rules before they are compared with the configured reference date, or the TeX
% compilation date when no fixed reference has been supplied, and normalized
% for display. Project-month values are classified only after a valid
% start date has been established. Project-month numbers are bounded before
% integer conversion. Canonical start-date storage separates project semantics
% from later display-format changes. Project-month boundaries use the start-day number where
% possible and the last valid day in shorter months, including the Gregorian
% leap-year rule. Every other value remains ordinary label text, which prevents
% descriptive or oversized labels from reaching numeric operations.
%
% Card surfaces, metadata chips, and deadline indicators are named
% \textsf{tcolorbox} styles. Seven empty, namespaced public styles are applied
% after the structural defaults, which permits controlled presentation changes
% without exposing internal environment names. The column-body style is used
% by the fixed renderer only; the other six styles are shared by both
% renderers. Semantic label text is copied
% into the same board snapshot as the other persistent configuration. Built-in
% header and semantic color pairs have been selected against a 4.5:1 reference
% contrast threshold. The clock symbol is a small inline TikZ drawing.
% Internal LaTeX2e control sequences consistently use the |kanban@| or |kb@|
% prefix, while LaTeX3 variables and functions use the |__kanban| module name.
% The public namespace is limited to the environment and commands documented in
% the user-interface sections. Nested boards are detected before local board
% state is initialized, preventing an inner environment from corrupting the
% active outer collection. The regression suite verifies configuration scope,
% diagnostics, canonical date snapshots, fixed reference dates, metadata flag
% semantics, malformed dimensions, compatibility with the declared LaTeX format,
% all three supported engines, all five date formats, explicit custom heading
% foregrounds, full-width cards, synchronized break-layout headings, actual
% cross-page continuation, and impossible-gap recovery. When Poppler is
% available, continuous integration also rasterizes the manual and the first
% page of every example and retains those images for visual review; exact pixel
% equality is intentionally not asserted across differing TeX and rasterizer
% releases. The release archive is inspected for its required source tests.
%
% \subsection{Bootstrap state}
%
% Bootstrap code establishes the conditions under which every later subsystem
% operates. It identifies the package, loads the packages on which rendering
% depends, and allocates the persistent and board-local state used by the
% configuration, collection, validation, and rendering stages. This subsection
% presents each new kind of definition separately. In every case, the
% functional responsibility is stated first, followed by the implementation
% method and the meaning of the LaTeX commands used to express it.
%
% \paragraph{Package declaration and runtime dependencies.}
% The first responsibility is to make the generated style file identifiable
% and to prevent it from running on a LaTeX format older than the declared
% compatibility boundary. The package must also load each facility that its
% public behavior requires rather than relying on another package to load that
% facility incidentally.
%
% The docstrip guard |%<*package>| begins the part of the documented source
% that is copied into |kanban.sty|; the corresponding closing guard appears at
% the end of the implementation. Within that guarded region,
% |\NeedsTeXFormat| declares the minimum LaTeX format date.
% |\ProvidesPackage| records the package name, release date, version, and short
% description. LaTeX prints this information in the compilation log and uses
% it when package versions are requested. Each |\RequirePackage| loads one
% declared dependency, while allowing LaTeX to avoid loading a dependency a
% second time if it is already present. Finally, |\tcbuselibrary| activates
% optional components of |tcolorbox| that are not part of its base interface.
%    \begin{macrocode}
%<*package>
\NeedsTeXFormat{LaTeX2e}[2020-10-01]
\ProvidesPackage{kanban}[2026/07/18 v1.0.0 Kanban project boards]
\RequirePackage{xcolor}
\RequirePackage{etoolbox}
\RequirePackage{tikz}
\RequirePackage{tcolorbox}
\RequirePackage{paracol}
\tcbuselibrary{skins,raster}
%    \end{macrocode}
%
% The dependencies have distinct responsibilities. |xcolor| provides named
% colors and color expressions. |etoolbox| supplies booleans, string tests,
% expandable list construction, and other LaTeX2e programming utilities.
% |tikz| draws the small clock symbol used by deadlines. |tcolorbox| constructs
% headers, cards, chips, and the fixed-board raster. |paracol| provides the
% independent, page-breakable column streams used by the |break| renderer.
% The |skins| library enables the required box decoration features, and the
% |raster| library provides the row-and-column layout for fixed boards.
%
% \paragraph{Built-in color constants.}
% The next definitions establish the package's semantic colors and the colors
% from which the built-in palettes are assembled. Defining them once gives the
% implementation stable symbolic names: later code refers to
% |kanbanHigh| or |kanbanSlate| instead of repeating numerical color values.
% This also keeps semantic meaning separate from the commands that draw a chip
% or heading.
%
% |\definecolor| is provided by |xcolor|. Its first argument creates the color
% name, the second selects the input model, and the third supplies the value in
% that model. The |HTML| model accepts six hexadecimal digits in red--green--
% blue order. These are ordinary named colors after their definitions, so they
% may be used by |tcolorbox|, TikZ, and standard LaTeX color commands.
%    \begin{macrocode}
\definecolor{kanbanInk}{HTML}{263238}
\definecolor{kanbanHigh}{HTML}{B23A48}
\definecolor{kanbanMed}{HTML}{E8A06A}
\definecolor{kanbanLow}{HTML}{477455}
\definecolor{kanbanSlate}{HTML}{4C5A6B}
\definecolor{kanbanBlue}{HTML}{3B7CA6}
\definecolor{kanbanTeal}{HTML}{267363}
\definecolor{kanbanSage}{HTML}{4F7D43}
\definecolor{kanbanAmber}{HTML}{876515}
\definecolor{kanbanClay}{HTML}{9B5030}
\definecolor{kanbanRose}{HTML}{AF5C74}
\definecolor{kanbanPlum}{HTML}{7C5C95}
\definecolor{kanbanMutedBlue}{HTML}{52758C}
\definecolor{kanbanMutedGreen}{HTML}{597C6E}
\definecolor{kanbanMutedPurple}{HTML}{746680}
\definecolor{kanbanMutedBrown}{HTML}{8A6F5A}
%    \end{macrocode}
%
% The three semantic colors represent overdue or high-priority work, work due
% at the active reference date, and completed or low-priority work. The other
% names are palette entries. Text colors are not hard-coded into the palette
% registry here because a later starred palette definition may provide an
% explicit foreground for each background.
%
% \paragraph{Boolean flags, board identity, and dimensions.}
% Rendering needs a small set of state flags. They record whether a board uses
% the page-breakable renderer, whether task declarations currently belong to
% an active board, whether the environment should render when it closes, and
% whether tasks inside a rejected nested board must be discarded. Two further
% flags carry the result of deadline classification. A unique board number
% prevents temporary data belonging to different boards from sharing names,
% and dedicated dimensions hold the persistent gap, its board-local snapshot,
% a validation candidate, and the combined width of all gaps.
%
% |\newbool| is an |etoolbox| definition command. It creates a named Boolean
% that is subsequently changed with |\booltrue| or |\boolfalse| and inspected
% with |\ifbool| or |\ifboolexpr|. |\newcounter| allocates a LaTeX counter;
% unlike a raw count register, it can be advanced through the standard counter
% interface. The internal macro |\kanban@boardgroup| stores the printable value
% associated with the active counter. |\newlength| allocates a skip register,
% and |\setlength| assigns the initial document-wide column gap. Separate
% registers are important here: validation can inspect a candidate without
% modifying a previously accepted value.
% \Needspace{16\baselineskip}
%    \begin{macrocode}
\newbool{kanban@break}
\newbool{kanban@active}
\newbool{kanban@render}
\newbool{kanban@discardtasks}
\newbool{kb@overdue}
\newbool{kb@duenow}
\newcounter{kanban@board}
\def\kanban@boardgroup{0}
\newlength\kanban@columnsep
\newlength\kanban@boardcolumnsep
\newlength\kanban@candidatecolumnsep
\newlength\kanban@totalgaps
\setlength\kanban@columnsep{2mm}
%    \end{macrocode}
%
% The |@| in these names marks them as internal LaTeX2e control sequences. It
% is a letter while a package file is read, but it is not normally a letter in
% a user's document. This convention reduces accidental use without presenting
% the names as a supported public API.
%
% \paragraph{Persistent defaults and board-local scalar snapshots.}
% Configuration values such as the start date, reference date, date format,
% visual mode, card colors, and semantic labels are simple token values. Most
% occur in pairs: one member is the persistent default for later boards, and
% the other is the snapshot used by the board currently being rendered. This
% separation implements the documented package--setup--board precedence and
% prevents a local override from altering subsequent boards.
%
% These internal scalar values are initialized with the TeX primitive |\def|.
% A definition such as |\def\kanban@dateformat{us}| creates a parameterless
% macro whose replacement text is |us|. Such macros are appropriate for small
% values that will be copied or reassigned frequently. Unlike |\newcommand|,
% |\def| does not check whether a command already exists; that is acceptable
% for private, package-owned names established in one controlled bootstrap
% location. Public commands are defined later with checking interfaces such as
% |\NewDocumentCommand|. Because the following definitions are executed at
% package load time, they establish document-wide initial values. Assignments
% made inside a board environment can remain local to that environment's TeX
% group.
%    \begin{macrocode}
\def\kanban@startdate{}
\def\kanban@thestart{}
\def\kanban@referencedate{}
\def\kanban@thereference{}
\def\kanban@dateformat{us}
\def\kanban@boarddateformat{us}
\def\kanban@mode{light}
\def\kanban@boardmode{light}
\def\kanban@textcolor{kanbanInk}
\def\kanban@cardcolor{white}
\def\kanban@bordercolor{black!15}
\def\kanban@authorback{black!6}
\def\kanban@authortext{black!58}
\def\kanban@deadlineback{black!7}
\def\kanban@deadlinetext{black!62}
\def\kanban@headingtext{white}
\def\kanban@hightext{white}
\def\kanban@mediumtext{kanbanInk}
\def\kanban@lowtext{white}
\def\kanban@highlabel{HIGH}
\def\kanban@mediumlabel{MED}
\def\kanban@lowlabel{LOW}
\def\kanban@donelabel{DONE}
\def\kanban@boardhighlabel{HIGH}
\def\kanban@boardmediumlabel{MED}
\def\kanban@boardlowlabel{LOW}
\def\kanban@boarddonelabel{DONE}
%    \end{macrocode}
%
% Empty date macros mean that no fixed date has been supplied. The default
% reference then comes from TeX's compilation date. The color macros contain
% either a named color or an |xcolor| expression and are expanded only when a
% box is rendered. The four persistent labels and their four board-local
% counterparts allow semantic wording to be customized without changing the
% priority and completion logic.
%
% \paragraph{LaTeX3 syntax and typed state containers.}
% Palettes and columns require structured data rather than isolated scalar
% macros. Palette names must map to lists of colors, columns must map to numeric
% identifiers, and parsing routines need temporary lists, token values,
% sequences, Boolean flags, and integers. The implementation therefore enters
% LaTeX3 syntax and allocates a typed variable for every distinct role.
%
% |\ExplSyntaxOn| changes the programming syntax so underscores and colons may
% appear in control-sequence names and ordinary spaces are ignored. LaTeX3
% names describe both scope and data type. The prefix |g__kanban| denotes
% private global package state, whereas |l__kanban| denotes private local or
% temporary state. These prefixes state the intended scope; the local or
% global form of the assignment command still determines the actual TeX
% assignment. The suffixes |_prop|, |_clist|, |_tl|, |_seq|, |_bool|, and
% |_int| identify a property list, comma list, token list, sequence, Boolean,
% and integer respectively. Function names use a colon followed by an argument
% signature; for example, |:nN| means one braced token-list argument followed
% by one control-sequence argument.
%
% Each allocation command creates one variable and reports an error if that
% variable already exists. |\prop_new:N| creates a key--value mapping;
% |\clist_new:N| creates a comma-separated list; |\tl_new:N| creates a general
% token-list variable; |\seq_new:N| creates an ordered sequence;
% |\bool_new:N| creates a LaTeX3 Boolean; and |\int_new:N| creates an integer.
% The final |N| in each command signature means that the argument must be the
% variable control sequence itself, not its textual name.
% \Needspace{30\baselineskip}
%    \begin{macrocode}
\ExplSyntaxOn
\prop_new:N \g__kanban_palettes_prop
\prop_new:N \g__kanban_palette_texts_prop
\clist_new:N \g__kanban_palette_clist
\clist_new:N \g__kanban_palette_fg_clist
\tl_new:N \g__kanban_palette_name_tl
\clist_new:N \l__kanban_board_palette_clist
\clist_new:N \l__kanban_board_palette_fg_clist
\clist_new:N \l__kanban_palette_bg_clist
\clist_new:N \l__kanban_palette_fg_clist
\seq_new:N \l__kanban_palette_pair_seq
\bool_new:N \l__kanban_bad_palette_bool
\clist_new:N \g__kanban_columns_clist
\clist_new:N \l__kanban_board_columns_clist
\prop_new:N \l__kanban_column_ids_prop
\int_new:N \l__kanban_column_index_int
\tl_new:N \l__kanban_palette_name_tl
\tl_new:N \l__kanban_column_tl
\tl_new:N \l__kanban_column_id_tl
\tl_new:N \l__kanban_choice_tl
\tl_new:N \l__kanban_pending_startdate_tl
\tl_new:N \l__kanban_pending_reference_date_tl
\bool_new:N \l__kanban_pending_startdate_bool
\bool_new:N \l__kanban_pending_reference_date_bool
%    \end{macrocode}
%
% The two global property lists form the palette registry: one stores heading
% backgrounds and the other stores optional heading foregrounds. Global comma
% lists hold the currently selected persistent palette and column defaults.
% Their local counterparts are the snapshot for an active board. The remaining
% local containers support palette-pair validation, column normalization,
% numeric column lookup, normalized choices, and the deferred processing of
% start and reference dates. Deferring those dates is necessary because a
% later |date-format| key in the same option list must be resolved before an
% ambiguous date is interpreted.
%
% \paragraph{The initial column model.}
% An environment with no |columns| key must still be valid. The package
% therefore installs the three canonical default stages as soon as the global
% column list has been allocated. The names are stored in normalized uppercase
% form so task assignment can be case-insensitive from the beginning.
%
% |\clist_gset:Nn| assigns a comma list globally. Its |N| argument is the
% destination variable, its |n| argument is literal braced input, and the |g|
% indicates that the assignment survives TeX grouping. Later setup commands
% use the same global list as the persistent source from which each board takes
% its local snapshot.
%    \begin{macrocode}
\clist_gset:Nn \g__kanban_columns_clist {BACKLOG,DOING,DONE}
%    \end{macrocode}
%
% \paragraph{Named diagnostics.}
% The bootstrap stage finally defines reusable diagnostics for errors that may
% arise in several configuration scopes. Naming each diagnostic keeps its main
% message and corrective help text consistent whether invalid input originates
% in a package option, |\kbset|, or a board option.
%
% |\msg_new:nnnn| creates a LaTeX3 message. The first argument is the module
% name, the second is the message identifier, the third is the diagnostic text,
% and the fourth is explanatory help. Parameter markers such as |#1| and |#2|
% are placeholders supplied when the message is issued. Under
% |\ExplSyntaxOn|, a tilde represents an explicit interword space. Defining a
% message does not print it; later code selects its severity with commands such
% as |\msg_warning:nnnn| and supplies the placeholder values.
% \Needspace{14\baselineskip}
%    \begin{macrocode}
\msg_new:nnnn {kanban}{invalid-date-setting}
  {Invalid~#1~'#2';~keeping~the~current~value.}
  {Use~a~valid~Gregorian~date~in~one~of~the~documented~numeric~forms.}
\msg_new:nnnn {kanban}{invalid-dimension}
  {Invalid~#1~value~'#2';~keeping~the~current~value.}
  {Use~a~non-negative~explicit~dimension,~such~as~2mm,~or~a~registered
    ~LaTeX~length.}
\msg_new:nnnn {kanban}{unknown-key}
  {Unknown~#1~option~'#2';~the~option~was~ignored.}
  {Check~the~spelling~and~the~documented~scope~of~the~key.}
%    \end{macrocode}
%
% At this point the package has allocated, but has not yet populated, the
% palette registries and board-local lookup structures. The following
% subsections define the functions that validate and populate those containers.
%
% \subsection{Palette and column registries}
%
% Palette names map to background lists and, for starred definitions, paired
% heading-foreground lists in global property lists. Column names are
% normalized separately and receive board-local numeric identifiers; this
% prevents user text from becoming part of an internal control-sequence name.
% The registry layer consequently has two responsibilities. It maintains the
% persistent palette catalogue, and it converts the textual column interface
% into a compact board-local representation suitable for deferred rendering.
%
% \paragraph{A common palette diagnostic.}
% Palette validation has several failure paths, but they should all identify
% the package in the same way. The small helper below centralizes that
% responsibility and leaves the more substantial functions concerned with
% validation and storage.
%
% |\cs_new_protected:Npn| defines a new protected LaTeX3 function. The
% signature |:n| states that the function accepts one braced token-list
% argument. The |Npn| part has the same broad shape as a protected TeX
% definition: |N| is the function being defined, |p| denotes its parameter
% text, and |n| denotes its unexpanded replacement text. Protection prevents
% the warning operation from expanding unexpectedly in an expansion-only
% context. |\PackageWarning| then writes a non-fatal package diagnostic to the
% terminal and log.
%    \begin{macrocode}
\cs_new_protected:Npn \__kanban_palette_warning:n #1
  {\PackageWarning{kanban}{#1}}
%    \end{macrocode}
%
% \paragraph{Plain palette definitions.}
% An unstarred palette definition must accept a name and a non-empty comma
% list of heading backgrounds. A successful definition replaces any palette
% of the same name, removes a previously associated foreground list, and, if
% the redefined palette is currently selected, refreshes the active persistent
% snapshot immediately. An invalid definition leaves all registries unchanged.
%
% The implementation first copies and trims the name with |\tl_set:Nn| and
% |\tl_trim_spaces:N|. |\tl_if_blank:VTF| examines the value held in the token
% list: |V| means ``use the value of this variable'', while |TF| supplies true
% and false branches. The color input is parsed by |\clist_set:Nn|, and
% |\clist_if_empty:NTF| distinguishes an empty comma list from a usable one.
% On success, |\prop_gput:NVV| globally associates the palette name with the
% color list; its |NVV| signature takes a property-list variable followed by
% the values of two variables. |\prop_gremove:NV| deletes stale foreground
% data. Finally, |\tl_if_eq:NNT| compares the new name with the selected name,
% |\clist_gset_eq:NN| copies the accepted list globally, and
% |\clist_gclear:N| records that the plain palette has no explicit heading
% foregrounds.
%    \begin{macrocode}
\cs_new_protected:Npn \__kanban_palette_define:nn #1#2
  {
    \tl_set:Nn \l__kanban_palette_name_tl {#1}
    \tl_trim_spaces:N \l__kanban_palette_name_tl
    \tl_if_blank:VTF \l__kanban_palette_name_tl
      {\__kanban_palette_warning:n {palette~name~must~not~be~empty}}
      {
        \clist_set:Nn \l_tmpa_clist {#2}
        \clist_if_empty:NTF \l_tmpa_clist
          {\__kanban_palette_warning:n
            {palette~'\tl_use:N \l__kanban_palette_name_tl'
              \c_space_tl must~contain~at~least~one~color}}
          {
            \prop_gput:NVV \g__kanban_palettes_prop
              \l__kanban_palette_name_tl \l_tmpa_clist
            \prop_gremove:NV \g__kanban_palette_texts_prop
              \l__kanban_palette_name_tl
            \tl_if_eq:NNT \l__kanban_palette_name_tl
              \g__kanban_palette_name_tl
              {
                \clist_gset_eq:NN \g__kanban_palette_clist
                  \l_tmpa_clist
                \clist_gclear:N \g__kanban_palette_fg_clist
              }
          }
      }
  }
%    \end{macrocode}
%
% \paragraph{Paired palette definitions.}
% The starred public syntax associates every heading background with an
% explicit heading foreground. Its validator therefore requires each comma
% item to contain exactly one |/| separator and two non-empty components. The
% complete definition is committed only after every pair has passed; a single
% malformed item rejects the entire palette rather than installing a partial
% result.
%
% The background and foreground working lists are emptied with
% |\clist_clear:N|, and |\bool_set_false:N| initializes the transaction's
% failure flag. The comma-list mapping command |\clist_map_inline:Nn| then
% visits each pair. Inside its body, |##1| is the current list item because the
% mapping code itself occurs within a macro definition.
% |\seq_set_split:Nnn| splits that item at |/|, while
% |\seq_count:N| and |\int_compare:nNnTF| require exactly two sequence
% elements. |\seq_item:Nn| retrieves each component. The predicate forms
% |\tl_if_blank_p:V| return Boolean expressions, and |\bool_if:nTF| combines
% them with the LaTeX3 logical disjunction operator. Accepted components are
% appended with |\clist_put_right:NV|. After the map,
% |\bool_if:NTF| either emits one diagnostic or globally commits both property
% entries and, where necessary, refreshes the selected palette lists.
%    \begin{macrocode}
\cs_new_protected:Npn \__kanban_palette_define_pairs:nn #1#2
  {
    \tl_set:Nn \l__kanban_palette_name_tl {#1}
    \tl_trim_spaces:N \l__kanban_palette_name_tl
    \clist_clear:N \l__kanban_palette_bg_clist
    \clist_clear:N \l__kanban_palette_fg_clist
    \bool_set_false:N \l__kanban_bad_palette_bool
    \tl_if_blank:VTF \l__kanban_palette_name_tl
      {\__kanban_palette_warning:n {palette~name~must~not~be~empty}}
      {
        \clist_set:Nn \l_tmpa_clist {#2}
        \clist_if_empty:NTF \l_tmpa_clist
          {\__kanban_palette_warning:n
            {palette~'\tl_use:N \l__kanban_palette_name_tl'
              \c_space_tl must~contain~at~least~one~color/text~pair}}
          {
            \clist_map_inline:Nn \l_tmpa_clist
              {
                \seq_set_split:Nnn \l__kanban_palette_pair_seq {/}{##1}
                \int_compare:nNnTF
                  {\seq_count:N \l__kanban_palette_pair_seq} = {2}
                  {
                    \tl_set:Nx \l_tmpa_tl
                      {\seq_item:Nn \l__kanban_palette_pair_seq {1}}
                    \tl_set:Nx \l_tmpb_tl
                      {\seq_item:Nn \l__kanban_palette_pair_seq {2}}
                    \tl_trim_spaces:N \l_tmpa_tl
                    \tl_trim_spaces:N \l_tmpb_tl
                    \bool_if:nTF
                      {
                        \tl_if_blank_p:V \l_tmpa_tl
                        || \tl_if_blank_p:V \l_tmpb_tl
                      }
                      {\bool_set_true:N \l__kanban_bad_palette_bool}
                      {
                        \clist_put_right:NV
                          \l__kanban_palette_bg_clist \l_tmpa_tl
                        \clist_put_right:NV
                          \l__kanban_palette_fg_clist \l_tmpb_tl
                      }
                  }
                  {\bool_set_true:N \l__kanban_bad_palette_bool}
              }
            \bool_if:NTF \l__kanban_bad_palette_bool
              {\__kanban_palette_warning:n
                {palette~'\tl_use:N \l__kanban_palette_name_tl'
                  \c_space_tl must~use~non-empty~color/text~pairs;
                  \c_space_tl definition~ignored}}
              {
                \prop_gput:NVV \g__kanban_palettes_prop
                  \l__kanban_palette_name_tl
                  \l__kanban_palette_bg_clist
                \prop_gput:NVV \g__kanban_palette_texts_prop
                  \l__kanban_palette_name_tl
                  \l__kanban_palette_fg_clist
                \tl_if_eq:NNT \l__kanban_palette_name_tl
                  \g__kanban_palette_name_tl
                  {
                    \clist_gset_eq:NN \g__kanban_palette_clist
                      \l__kanban_palette_bg_clist
                    \clist_gset_eq:NN \g__kanban_palette_fg_clist
                      \l__kanban_palette_fg_clist
                  }
              }
          }
      }
  }
%    \end{macrocode}
%
% \paragraph{Persistent selection and board-local selection.}
% Selecting a palette is distinct from defining one. Document-wide setup must
% replace the persistent defaults used by future boards, whereas a board option
% must replace only the current board's local snapshot. Both operations verify
% the requested name against the same global registry and preserve the current
% selection when the name is unknown.
%
% |\prop_get:NVNTF| looks up the key held in a token-list variable. If the key
% exists, it writes the associated value to the supplied token-list variable
% and executes the true branch; otherwise it executes the false branch.
% Commands containing |gset| or |gclear| make global assignments for persistent
% configuration. Their local |set| and |clear| counterparts are used for the
% active board. |\tl_gset_eq:NN| records the selected name by equivalence, and
% |\clist_set_eq:NN| in |\kanban@paletteboardinit| takes the initial local
% snapshot when a board begins.
%    \begin{macrocode}
\cs_new_protected:Npn \__kanban_palette_select_global:n #1
  {
    \tl_set:Nn \l__kanban_palette_name_tl {#1}
    \tl_trim_spaces:N \l__kanban_palette_name_tl
    \prop_get:NVNTF \g__kanban_palettes_prop
      \l__kanban_palette_name_tl \l_tmpa_tl
      {
        \clist_gset:NV \g__kanban_palette_clist \l_tmpa_tl
        \prop_get:NVNTF \g__kanban_palette_texts_prop
          \l__kanban_palette_name_tl \l_tmpb_tl
          {\clist_gset:NV \g__kanban_palette_fg_clist \l_tmpb_tl}
          {\clist_gclear:N \g__kanban_palette_fg_clist}
        \tl_gset_eq:NN \g__kanban_palette_name_tl
          \l__kanban_palette_name_tl
      }
      {\__kanban_palette_warning:n
        {unknown~palette~'\tl_use:N \l__kanban_palette_name_tl';
          \c_space_tl keeping~the~current~palette}}
  }
\cs_new_protected:Npn \__kanban_palette_select_board:n #1
  {
    \tl_set:Nn \l__kanban_palette_name_tl {#1}
    \tl_trim_spaces:N \l__kanban_palette_name_tl
    \prop_get:NVNTF \g__kanban_palettes_prop
      \l__kanban_palette_name_tl \l_tmpa_tl
      {
        \clist_set:NV \l__kanban_board_palette_clist \l_tmpa_tl
        \prop_get:NVNTF \g__kanban_palette_texts_prop
          \l__kanban_palette_name_tl \l_tmpb_tl
          {\clist_set:NV \l__kanban_board_palette_fg_clist \l_tmpb_tl}
          {\clist_clear:N \l__kanban_board_palette_fg_clist}
      }
      {\__kanban_palette_warning:n
        {unknown~palette~'\tl_use:N \l__kanban_palette_name_tl';
          \c_space_tl keeping~the~current~palette}}
  }
\cs_new_protected:Npn \kanban@paletteboardinit
  {
    \clist_set_eq:NN \l__kanban_board_palette_clist
      \g__kanban_palette_clist
    \clist_set_eq:NN \l__kanban_board_palette_fg_clist
      \g__kanban_palette_fg_clist
  }
%    \end{macrocode}
%
% \paragraph{Cyclic palette lookup.}
% A board may contain more columns than its palette contains colors. The lookup
% function therefore maps each one-based column number cyclically onto the
% board palette and returns both the background and its heading foreground.
% If the selected palette has no explicit foreground list, the board theme's
% standard heading color is used.
%
% |\clist_count:N| obtains the palette length. |\int_mod:nn| performs the
% zero-based wraparound, |\int_eval:n| evaluates the resulting one-based index,
% and |\clist_item:Nn| retrieves the corresponding entry. The |x| assignment
% in |\tl_set:Nx| expands this calculation before storing it. |\cs_set:Npx|
% locally defines the LaTeX2e-facing bridge macros with an expanded replacement
% value. This bridge lets later |etoolbox| and |tcolorbox| code consume data
% prepared in the LaTeX3 registry without knowing its container types.
%    \begin{macrocode}
\cs_new_protected:Npn \kanban@palettecolor #1
  {
    \int_set:Nn \l_tmpa_int
      {\clist_count:N \l__kanban_board_palette_clist}
    \tl_set:Nx \l_tmpa_tl
      {
        \clist_item:Nn \l__kanban_board_palette_clist
          {\int_eval:n {1 + \int_mod:nn {#1 - 1} {\l_tmpa_int}}}
      }
    \cs_set:Npx \kb@thiscolor {\tl_use:N \l_tmpa_tl}
    \clist_if_empty:NTF \l__kanban_board_palette_fg_clist
      {\cs_set:Npn \kb@thisheadingtext {\kanban@headingtext}}
      {
        \tl_set:Nx \l_tmpa_tl
          {
            \clist_item:Nn \l__kanban_board_palette_fg_clist
              {\int_eval:n {1 + \int_mod:nn {#1 - 1} {\l_tmpa_int}}}
          }
        \cs_set:Npx \kb@thisheadingtext {\tl_use:N \l_tmpa_tl}
      }
  }
%    \end{macrocode}
% \begin{macro}{\kanbandefinepalette}
% \begin{macro}{\kbdefinepalette}
% The public definition command exposes both registry formats through one
% interface. Its unstarred form accepts a list of backgrounds; its starred
% form accepts background/foreground pairs. The long command is the stable
% public name, while the concise alias is installed only when that name is
% free so that loading |kanban| does not overwrite another package's command.
%
% |\NewDocumentCommand| checks that the canonical command is new and parses
% the xparse signature |smm|: |s| records an optional star as a Boolean token,
% and the two |m| entries read mandatory arguments. |\IfBooleanTF| selects the
% paired or plain validator. |\cs_if_exist:NTF| tests whether the alias already
% exists; the occupied case uses |\PackageWarning| and preserves it, while
% |\cs_new_eq:NN| gives the free alias exactly the same meaning as the
% canonical command without creating a second implementation.
%    \begin{macrocode}
\NewDocumentCommand \kanbandefinepalette {smm}
  {
    \IfBooleanTF{#1}
      {\__kanban_palette_define_pairs:nn {#2}{#3}}
      {\__kanban_palette_define:nn {#2}{#3}}
  }
\cs_if_exist:NTF \kbdefinepalette
  {\PackageWarning{kanban}
    {\string\kbdefinepalette\space is~already~defined;
      \MessageBreak use~\string\kanbandefinepalette\space instead}}
  {\cs_new_eq:NN \kbdefinepalette \kanbandefinepalette}
%    \end{macrocode}
% \end{macro}
% \end{macro}
%
% \paragraph{Internal selection bridges.}
% The setup and board-key layers are written partly in traditional LaTeX2e
% syntax. These two protected wrappers provide stable internal names through
% which those layers invoke the LaTeX3 selection functions. They add no state
% or validation of their own.
%
% Both definitions again use |\cs_new_protected:Npn| with one braced argument.
% Passing |{#1}| to the underlying |:n| function deliberately preserves the
% argument as a token list until that function performs its own trimming and
% lookup.
%    \begin{macrocode}
\cs_new_protected:Npn \kanban@selectpalette #1
  {\__kanban_palette_select_global:n {#1}}
\cs_new_protected:Npn \kanban@selectboardpalette #1
  {\__kanban_palette_select_board:n {#1}}
%    \end{macrocode}
%
% \paragraph{Column-name normalization.}
% Column identity is intentionally case-insensitive. A single-name function
% removes surrounding spaces and converts the result to uppercase. A list
% function applies that rule to every item, ignores empty items, and diagnoses
% duplicates after normalization. Thus |Doing| and |DOING| cannot silently
% create two visually different names that refer to the same logical column.
%
% The signature |:nN| means that the first function receives textual input and
% a destination variable. |\tl_upper_case:n| performs the Unicode-aware
% LaTeX3 case conversion; |\exp_args:NV| supplies the current value of the
% destination as its argument, and |\tl_set:Nx| stores the expanded result.
% For lists, |\clist_map_inline:nn| maps directly over literal input,
% |\tl_if_blank:VF| skips blank normalized values, and
% |\clist_if_in:NVTF| checks the accumulating destination. A duplicate is a
% user error rather than a recoverable preference, so |\PackageError| provides
% both a concise error and corrective help. Unique values enter the list via
% |\clist_put_right:NV|.
%    \begin{macrocode}
\cs_new_protected:Npn \__kanban_column_normalize:nN #1#2
  {
    \tl_set:Nn #2 {#1}
    \tl_trim_spaces:N #2
    \tl_set:Nx #2 {\exp_args:NV \tl_upper_case:n #2}
  }
\cs_new_protected:Npn \__kanban_columns_normalize:nN #1#2
  {
    \clist_clear:N #2
    \clist_map_inline:nn {#1}
      {
        \__kanban_column_normalize:nN {##1} \l__kanban_column_tl
        \tl_if_blank:VF \l__kanban_column_tl
          {
            \clist_if_in:NVTF #2 \l__kanban_column_tl
              {\PackageError{kanban}
                {duplicate~kanban~column~
                  '\tl_use:N \l__kanban_column_tl'}
                {Column~names~are~case-insensitive.\MessageBreak
                  Remove~the~repeated~name~from~the~column~list.}}
              {\clist_put_right:NV #2 \l__kanban_column_tl}
          }
      }
  }
%    \end{macrocode}
%
% \paragraph{Persistent and board-local column lists.}
% As with palettes, columns have a persistent default and a board-local
% snapshot. Document setup accepts only a normalized non-empty list and makes
% it global. Board initialization copies that default locally; a non-blank
% board override may then replace the copy without affecting later boards.
%
% |\clist_gset:NV| is the global value-copying assignment used after persistent
% validation. |\clist_set_eq:NN| copies a list locally for the active board.
% |\tl_if_blank:nF| and |\clist_if_empty:NF| are one-branch conditionals: their
% bodies run only when the literal argument or normalized list, respectively,
% is not blank or empty. The short |\kanban@selectcolumns| wrapper exposes the
% global selector to the setup key layer.
%    \begin{macrocode}
\cs_new_protected:Npn \__kanban_columns_select_global:n #1
  {
    \__kanban_columns_normalize:nN {#1} \l_tmpa_clist
    \clist_if_empty:NTF \l_tmpa_clist
      {\PackageWarning{kanban}
        {column~list~must~contain~at~least~one~name;
          ~keeping~the~current~list}}
      {\clist_gset:NV \g__kanban_columns_clist \l_tmpa_clist}
  }
\cs_new_protected:Npn \kanban@selectcolumns #1
  {\__kanban_columns_select_global:n {#1}}
\cs_new_protected:Npn \kanban@columnsboardinit
  {\clist_set_eq:NN \l__kanban_board_columns_clist
    \g__kanban_columns_clist}
\cs_new_protected:Npn \kanban@selectboardcolumns #1
  {
    \tl_if_blank:nF {#1}
      {
        \__kanban_columns_normalize:nN {#1} \l_tmpa_clist
        \clist_if_empty:NF \l_tmpa_clist
          {\clist_set_eq:NN
            \l__kanban_board_columns_clist \l_tmpa_clist}
      }
  }
%    \end{macrocode}
%
% \paragraph{Task-target resolution.}
% A task names a column textually, but collection must append the task to a
% numeric board-local list. The resolver normalizes the supplied name, exposes
% that normalized text for diagnostics, and looks it up in the active board's
% name-to-identifier property list. A missing key leaves the identifier empty,
% allowing the collection layer to issue the context-specific error.
%
% |\prop_get:NVNTF| performs the lookup. A successful value is written to the
% token-list variable |\l__kanban_column_id_tl|. |\cs_set:Npn| first
% establishes an empty result;
% |\cs_set:Npx| then expands and copies a successful normalized name or numeric
% identifier into the LaTeX2e-facing task macros. These assignments are local
% to the surrounding board group.
%    \begin{macrocode}
\cs_new_protected:Npn \kanban@normtaskcol #1
  {
    \__kanban_column_normalize:nN {#1} \l__kanban_column_tl
    \cs_set:Npx \kbtask@column {\tl_use:N \l__kanban_column_tl}
    \cs_set:Npn \kbtask@columnid { }
    \prop_get:NVNTF
      \l__kanban_column_ids_prop \l__kanban_column_tl
      \l__kanban_column_id_tl
      {\cs_set:Npx \kbtask@columnid
        {\tl_use:N \l__kanban_column_id_tl}}
      { }
  }
%    \end{macrocode}
%
% \paragraph{Numeric column registration.}
% Before tasks can be collected, every normalized board column receives a
% consecutive identifier. The setup function also prepares a printable list of
% valid names for error messages and delegates the creation of per-column
% counters, lists, colors, and renderer entries to |\kanban@setupcolumn|, which
% is defined after the board key family.
%
% |\clist_use:Nn| serializes the board list with a comma delimiter.
% |\prop_clear:N| and |\int_zero:N| reset the local lookup and counter.
% |\clist_map_inline:Nn| iterates over the stored list, |\int_incr:N| advances
% the identifier, and |\prop_put:Nnx| stores the expanded integer value under
% the current normalized name. |\int_use:N| converts the LaTeX3 integer to
% character tokens. Finally, |\exp_args:Nnx| expands only the numeric argument
% before calling the later LaTeX2e setup function.
%    \begin{macrocode}
\cs_new_protected:Npn \kanban@setupboardcolumns
  {
    \cs_set:Npx \kanban@availablecolumns
      {\clist_use:Nn \l__kanban_board_columns_clist {,}}
    \prop_clear:N \l__kanban_column_ids_prop
    \int_zero:N \l__kanban_column_index_int
    \clist_map_inline:Nn \l__kanban_board_columns_clist
      {
        \int_incr:N \l__kanban_column_index_int
        \prop_put:Nnx \l__kanban_column_ids_prop {##1}
          {\int_use:N \l__kanban_column_index_int}
        \exp_args:Nnx \kanban@setupcolumn {##1}
          {\int_use:N \l__kanban_column_index_int}
      }
  }
%    \end{macrocode}
%
% \paragraph{Choice normalization and registry cleanup.}
% Several key values, including modes and date-format names, share the same
% case-insensitive choice syntax. The choice helper trims and lowercases such a
% value before publishing it as |\kb@choice|. At the other end of the board
% lifecycle, the cleanup function removes all dynamically created per-column
% control sequences so that completed boards do not retain card material or
% counters in the global hash table.
%
% |\tl_lower_case:n| is applied through the same value-expansion pattern used
% for column uppercasing. |\int_step_inline:nn| executes its body once for each
% integer from one through |\kanban@ncols|; its current value is written |##1|
% inside this macro definition. |\cs_undefine:c| removes a control sequence
% whose name is constructed from character tokens. The |c| argument type is
% essential because the board number and column identifier are data, not
% literal control-sequence tokens.
%    \begin{macrocode}
\cs_new_protected:Npn \kanban@normalizechoice #1
  {
    \tl_set:Nn \l__kanban_choice_tl {#1}
    \tl_trim_spaces:N \l__kanban_choice_tl
    \tl_set:Nx \l__kanban_choice_tl
      {\exp_args:NV \tl_lower_case:n \l__kanban_choice_tl}
    \cs_set:Npx \kb@choice {\tl_use:N \l__kanban_choice_tl}
  }
\cs_new_protected:Npn \kanban@cleanupboard
  {
    \int_step_inline:nn {\kanban@ncols}
      {
        \cs_undefine:c {kanban@name@##1}
        \cs_undefine:c {kanban@color@##1}
        \cs_undefine:c {kanban@headingtext@##1}
        \cs_undefine:c
          {kanban@list@\kanban@boardgroup @##1}
        \cs_undefine:c
          {kanban@count@\kanban@boardgroup @##1}
      }
  }
%    \end{macrocode}
%
% \subsection{Deadline engine}
%
% Deadline classification uses exact regular expressions before any numeric
% conversion. The five supported notations enter one component validator and
% one formatter; bounded project-month labels use a separate validated path,
% while arbitrary descriptive labels remain neutral and unchanged. The engine
% deliberately separates syntax recognition, Gregorian validation, display
% formatting, and temporal classification. This separation permits a task to
% retain a neutral descriptive deadline without misclassifying it as a broken
% date.
%
% \paragraph{Typed parsing and calendar state.}
% Date processing requires the original token list, regular-expression capture
% groups, numeric components, maximum-day calculations, project-month state,
% reference-date state, and explicit status flags. Allocating a distinct typed
% variable for each role makes intermediate state visible and prevents the
% same scratch register from carrying unrelated meanings across nested tests.
%
% |\tl_new:N| allocates the normalized input token list. |\seq_new:N| creates
% the sequence into which a regular expression writes its complete match and
% capture groups. Each |\int_new:N| allocates an integer variable for a
% calendar or project component. Each |\bool_new:N| allocates a LaTeX3 Boolean
% used to distinguish valid data from recognized-but-invalid syntax and from
% neutral text. Allocation initializes these variables to empty, zero, or
% false according to their type; individual parsing operations still reset the
% relevant flags explicitly.
%    \begin{macrocode}
\tl_new:N \l__kanban_deadline_tl
\seq_new:N \l__kanban_deadline_match_seq
\int_new:N \l__kanban_deadline_day_int
\int_new:N \l__kanban_deadline_month_int
\int_new:N \l__kanban_deadline_year_int
\int_new:N \l__kanban_deadline_max_day_int
\int_new:N \l__kanban_project_month_int
\int_new:N \l__kanban_project_start_day_int
\int_new:N \l__kanban_project_start_month_int
\int_new:N \l__kanban_project_start_year_int
\int_new:N \l__kanban_project_boundary_day_int
\int_new:N \l__kanban_reference_day_int
\int_new:N \l__kanban_reference_month_int
\int_new:N \l__kanban_reference_year_int
\bool_new:N \l__kanban_valid_date_bool
\bool_new:N \l__kanban_valid_project_month_bool
\bool_new:N \l__kanban_date_syntax_bool
\bool_new:N \l__kanban_ambiguous_slash_bool
\bool_new:N \l__kanban_project_month_syntax_bool
%    \end{macrocode}
%
% \paragraph{Maximum days in a Gregorian month.}
% Component validation needs the number of permitted days in the parsed month.
% The helper implements the Gregorian leap-year rule: February has 29 days in
% years divisible by 400 or divisible by 4 but not by 100; April, June,
% September, and November have 30; all remaining valid months have 31.
%
% |\int_case:nnF| compares its first integer expression with each listed case
% and executes the fallback |F| branch when no case matches. Nested
% |\int_compare:nNnTF| functions express the leap-year decisions, and
% |\int_mod:nn| supplies divisibility remainders. |\int_set:Nn| assigns the
% selected maximum to the dedicated integer. Braced integer expressions are
% evaluated by the LaTeX3 integer layer, so no explicit |\numexpr| is required.
%    \begin{macrocode}
\cs_new_protected:Npn \__kanban_set_max_day:
  {
    \int_case:nnF {\l__kanban_deadline_month_int}
      {
        {2}
          {
            \int_compare:nNnTF
              {
                \int_mod:nn
                  {\l__kanban_deadline_year_int}
                  {400}
              }
              = {0}
              {
                \int_set:Nn \l__kanban_deadline_max_day_int
                  {29}
              }
              {
                \int_compare:nNnTF
                  {
                    \int_mod:nn
                      {\l__kanban_deadline_year_int}
                      {100}
                  }
                  = {0}
                  {
                    \int_set:Nn \l__kanban_deadline_max_day_int
                      {28}
                  }
                  {
                    \int_compare:nNnTF
                      {
                        \int_mod:nn
                          {\l__kanban_deadline_year_int}
                          {4}
                      }
                      = {0}
                      {
                        \int_set:Nn \l__kanban_deadline_max_day_int
                          {29}
                      }
                      {
                        \int_set:Nn \l__kanban_deadline_max_day_int
                          {28}
                      }
                  }
              }
          }
        {4}{\int_set:Nn \l__kanban_deadline_max_day_int {30}}
        {6}{\int_set:Nn \l__kanban_deadline_max_day_int {30}}
        {9}{\int_set:Nn \l__kanban_deadline_max_day_int {30}}
        {11}{\int_set:Nn \l__kanban_deadline_max_day_int {30}}
      }
      {\int_set:Nn \l__kanban_deadline_max_day_int {31}}
  }
%    \end{macrocode}
%
% \paragraph{Gregorian component validation and capture conversion.}
% Once a notation has assigned year, month, and day, the validator checks that
% the year is positive, the month lies between 1 and 12, and the day lies in
% the range permitted by the month and leap-year calculation. The capture
% helper provides the single conversion path from regular-expression text to a
% LaTeX3 integer.
%
% |\bool_set_false:N| establishes failure as the initial result. The one-branch
% form |\int_compare:nNnT| executes its body only when the comparison is true;
% nesting these tests means that the maximum-day helper is never called with an
% invalid month. The final test sets the validity Boolean with
% |\bool_set_true:N|. In the capture helper, |\seq_item:Nn| retrieves a capture
% group and |\int_set:Nn| parses its digit tokens as an integer. Its |:Nn|
% signature accepts the destination integer variable followed by the
% one-based sequence position.
%    \begin{macrocode}
\cs_new_protected:Npn \__kanban_validate_date:
  {
    \bool_set_false:N \l__kanban_valid_date_bool
    \int_compare:nNnT {\l__kanban_deadline_year_int} > {0}
      {
        \int_compare:nNnT {\l__kanban_deadline_month_int} > {0}
          {
            \int_compare:nNnT {\l__kanban_deadline_month_int} < {13}
              {
                \__kanban_set_max_day:
                \int_compare:nNnT {\l__kanban_deadline_day_int} > {0}
                  {
                    \int_compare:nNnT
                      {\l__kanban_deadline_day_int}
                      < {\l__kanban_deadline_max_day_int + 1}
                      {\bool_set_true:N \l__kanban_valid_date_bool}
                  }
              }
          }
      }
  }
\cs_new_protected:Npn \__kanban_set_match:Nn #1#2
  {
    \int_set:Nn #1
      {\seq_item:Nn \l__kanban_deadline_match_seq {#2}}
  }
%    \end{macrocode}
%
% \paragraph{Recognition of the five date notations.}
% The central parser recognizes dotted day-first dates, hyphenated ISO dates,
% slash-separated year-first dates, and the two locale-dependent slash forms.
% It first records whether the input has date syntax, then assigns all three
% components, and only afterward invokes Gregorian validation. For an ambiguous
% slash date, the selected date format determines the first interpretation. If
% that interpretation is impossible, the parser swaps day and month once and
% validates the alternative. This fallback accepts unambiguous human input
% such as |31/12/2026| without weakening the documented display format.
%
% |\regex_extract_once:nnNTF| matches an anchored regular expression once and
% stores the full match plus parenthesized captures in a sequence. The
% |\exp_args:NnV| prefix passes the value of the normalized input variable to
% the regular-expression function. |\A| and |\Z| anchor the complete string;
% bracket expressions and repetition counts describe the numeric fields.
% Nested |TF| forms implement an ordered grammar without numeric conversion of
% unmatched text. |\str_if_eq_p:ee| supplies expandable equality predicates to
% |\bool_if:nTF|, while |\bool_if:nT| combines the ambiguity and validity flags.
% |\int_set_eq:NN| copies integer values during the one-time day/month swap.
%    \begin{macrocode}
\cs_new_protected:Npn \__kanban_parse_date:nn #1#2
  {
    \bool_set_false:N \l__kanban_valid_date_bool
    \bool_set_false:N \l__kanban_date_syntax_bool
    \bool_set_false:N \l__kanban_ambiguous_slash_bool
    \tl_set:Nx \l__kanban_deadline_tl {#1}
    \tl_trim_spaces:N \l__kanban_deadline_tl
    \exp_args:NnV \regex_extract_once:nnNTF
      {\A([0-9]{1,2})\.([0-9]{1,2})\.([0-9]{1,4})\Z}
      \l__kanban_deadline_tl
      \l__kanban_deadline_match_seq
      {
        \bool_set_true:N \l__kanban_date_syntax_bool
        \__kanban_set_match:Nn \l__kanban_deadline_day_int {2}
        \__kanban_set_match:Nn \l__kanban_deadline_month_int {3}
        \__kanban_set_match:Nn \l__kanban_deadline_year_int {4}
      }
      {
        \exp_args:NnV \regex_extract_once:nnNTF
          {\A([0-9]{4})-([0-9]{2})-([0-9]{2})\Z}
          \l__kanban_deadline_tl
          \l__kanban_deadline_match_seq
          {
            \bool_set_true:N \l__kanban_date_syntax_bool
            \__kanban_set_match:Nn \l__kanban_deadline_year_int {2}
            \__kanban_set_match:Nn \l__kanban_deadline_month_int {3}
            \__kanban_set_match:Nn \l__kanban_deadline_day_int {4}
          }
          {
            \exp_args:NnV \regex_extract_once:nnNTF
              {\A([0-9]{4})/([0-9]{2})/([0-9]{2})\Z}
              \l__kanban_deadline_tl
              \l__kanban_deadline_match_seq
              {
                \bool_set_true:N \l__kanban_date_syntax_bool
                \__kanban_set_match:Nn
                  \l__kanban_deadline_year_int {2}
                \__kanban_set_match:Nn
                  \l__kanban_deadline_month_int {3}
                \__kanban_set_match:Nn
                  \l__kanban_deadline_day_int {4}
              }
              {
                \exp_args:NnV \regex_extract_once:nnNT
                  {\A([0-9]{1,2})/([0-9]{1,2})/([0-9]{1,4})\Z}
                  \l__kanban_deadline_tl
                  \l__kanban_deadline_match_seq
                  {
                    \bool_set_true:N \l__kanban_date_syntax_bool
                    \bool_set_true:N \l__kanban_ambiguous_slash_bool
                    \__kanban_set_match:Nn
                      \l__kanban_deadline_year_int {4}
                    \bool_if:nTF
                      {
                        \str_if_eq_p:ee {#2} {english}
                        || \str_if_eq_p:ee {#2} {german}
                      }
                      {
                        \__kanban_set_match:Nn
                          \l__kanban_deadline_day_int {2}
                        \__kanban_set_match:Nn
                          \l__kanban_deadline_month_int {3}
                      }
                      {
                        \__kanban_set_match:Nn
                          \l__kanban_deadline_month_int {2}
                        \__kanban_set_match:Nn
                          \l__kanban_deadline_day_int {3}
                      }
                  }
              }
          }
      }
    \bool_if:NT \l__kanban_date_syntax_bool
      {\__kanban_validate_date:}
    \bool_if:nT
      {
        \l__kanban_ambiguous_slash_bool
        && ! \l__kanban_valid_date_bool
      }
      {
        \int_set_eq:NN \l_tmpa_int \l__kanban_deadline_day_int
        \int_set_eq:NN \l__kanban_deadline_day_int
          \l__kanban_deadline_month_int
        \int_set_eq:NN \l__kanban_deadline_month_int \l_tmpa_int
        \__kanban_validate_date:
      }
  }
%    \end{macrocode}
%
% \paragraph{Canonical and display formatting.}
% Internally, accepted dates are stored in one unambiguous year--month--day
% representation. At rendering time, the same numeric components are formatted
% according to the board's selected notation. Keeping canonical storage and
% presentation separate ensures that changing |date-format| affects input
% interpretation and display, but not chronological comparison.
%
% The two padding helpers are deliberately expandable functions defined with
% |\cs_new:Npn| rather than protected functions. |\int_compare:nNnT| emits a
% leading zero when required, and |\int_to_arabic:n| emits the decimal digits.
% The canonical function composes these expandable results directly. By
% contrast, |\kanban@preparedeadline| is protected because it performs
% assignments and parsing. It initially preserves the literal deadline, so
% invalid dates and descriptive labels remain printable. For a valid date,
% |\tl_set:Nx| materializes padded components, |\str_case:VnF| selects a format
% by the value of the board-format token list, and |\cs_set:Npx| defines the
% final display macro with the selected separators already expanded.
%    \begin{macrocode}
\cs_new:Npn \__kanban_two_digits:n #1
  {
    \int_compare:nNnT {#1} < {10} {0}
    \int_to_arabic:n {#1}
  }
\cs_new:Npn \__kanban_four_digits:n #1
  {
    \int_compare:nNnT {#1} < {1000} {0}
    \int_compare:nNnT {#1} < {100} {0}
    \int_compare:nNnT {#1} < {10} {0}
    \int_to_arabic:n {#1}
  }
\cs_new:Npn \__kanban_canonical_date:
  {
    \__kanban_four_digits:n {\l__kanban_deadline_year_int}
    -\__kanban_two_digits:n {\l__kanban_deadline_month_int}
    -\__kanban_two_digits:n {\l__kanban_deadline_day_int}
  }
\cs_new_protected:Npn \kanban@preparedeadline #1
  {
    \cs_set:Npn \kb@displaydeadline {#1}
    \__kanban_parse_date:nn {#1} {\kanban@boarddateformat}
    \bool_if:NT \l__kanban_valid_date_bool
      {
        \tl_set:Nx \l_tmpa_tl
          {\__kanban_two_digits:n {\l__kanban_deadline_day_int}}
        \tl_set:Nx \l_tmpb_tl
          {\__kanban_two_digits:n {\l__kanban_deadline_month_int}}
        \tl_set:Nx \l_tmpc_tl
          {\__kanban_four_digits:n {\l__kanban_deadline_year_int}}
        \tl_set:Nx \l_tmpd_tl {\kanban@boarddateformat}
        \str_case:VnF \l_tmpd_tl
          {
            {us}
              {\cs_set:Npx \kb@displaydeadline
                {\l_tmpb_tl/\l_tmpa_tl/\l_tmpc_tl}}
            {english}
              {\cs_set:Npx \kb@displaydeadline
                {\l_tmpa_tl/\l_tmpb_tl/\l_tmpc_tl}}
            {german}
              {\cs_set:Npx \kb@displaydeadline
                {\l_tmpa_tl.\l_tmpb_tl.\l_tmpc_tl}}
            {iso}
              {\cs_set:Npx \kb@displaydeadline
                {\l_tmpc_tl-\l_tmpb_tl-\l_tmpa_tl}}
            {ymd}
              {\cs_set:Npx \kb@displaydeadline
                {\l_tmpc_tl/\l_tmpb_tl/\l_tmpa_tl}}
          }
          { }
      }
  }
%    \end{macrocode}
%
% \paragraph{Project-month syntax.}
% Project deadlines such as |M1| are relative to a configured project start
% date rather than to a calendar date written on the card. The parser accepts
% a positive integer of at most six digits, records recognized but invalid
% |M|-prefixed numeric syntax separately, and leaves all other text neutral.
% The upper bound prevents pathological integer input while remaining far
% beyond ordinary project durations.
%
% |\regex_match:nnT| records the broad |M| followed by digits syntax, whereas
% |\regex_extract_once:nnNTF| applies the stricter nonzero bounded grammar and
% captures the integer. |\AM[0-9]+\Z| uses |+| for one or more digits;
% |[1-9][0-9]{0,5}| excludes zero and limits the accepted width. A successful
% capture is converted with |\seq_item:Nn| and |\int_set:Nn|, after which
% |\bool_set_true:N| marks the value as suitable for classification.
%    \begin{macrocode}
\cs_new_protected:Npn \__kanban_parse_project_month:n #1
  {
    \bool_set_false:N \l__kanban_valid_project_month_bool
    \bool_set_false:N \l__kanban_project_month_syntax_bool
    \tl_set:Nx \l__kanban_deadline_tl {#1}
    \tl_trim_spaces:N \l__kanban_deadline_tl
    \exp_args:NnV \regex_match:nnT
      {\AM[0-9]+\Z}
      \l__kanban_deadline_tl
      {\bool_set_true:N \l__kanban_project_month_syntax_bool}
    \exp_args:NnV \regex_extract_once:nnNTF
      {\AM([1-9][0-9]{0,5})\Z}
      \l__kanban_deadline_tl
      \l__kanban_deadline_match_seq
      {
        \int_set:Nn \l__kanban_project_month_int
          {\seq_item:Nn \l__kanban_deadline_match_seq {2}}
        \bool_set_true:N \l__kanban_valid_project_month_bool
      }
      { }
  }
%    \end{macrocode}
%
% \paragraph{Reference-date loading and calendar classification.}
% Classification is evaluated against a board-local reference date. If none
% was configured, the engine reads TeX's compilation year, month, and day. A
% configured value is already canonical, so the loader may extract it with one
% exact expression. Calendar dates are then converted to order-preserving
% integers of the form $10000y+100m+d$: a smaller value is overdue, and an
% equal value is due now.
%
% |\tl_if_blank:VTF| selects the live-date or configured-date path. |\year|,
% |\month|, and |\day| are TeX integer parameters.
%
% |\regex_extract_once:nnN| is used only in the configured branch. It is safe
% there because the setter has already validated and canonicalized every
% non-empty reference date. The classifier
% uses |\int_set:Nn| for both sortable values and nested
% |\int_compare:nNnTF| tests for the three possible temporal relations.
% |\booltrue| sets the existing |etoolbox| status flags consumed by the card
% renderer; leaving both false denotes a future or neutral deadline.
%    \begin{macrocode}
\cs_new_protected:Npn \__kanban_load_reference_date:
  {
    \tl_set:Nx \l_tmpa_tl {\kanban@thereference}
    \tl_if_blank:VTF \l_tmpa_tl
      {
        \int_set:Nn \l__kanban_reference_year_int {\year}
        \int_set:Nn \l__kanban_reference_month_int {\month}
        \int_set:Nn \l__kanban_reference_day_int {\day}
      }
      {
        \exp_args:NnV \regex_extract_once:nnN
          {\A([0-9]{4})-([0-9]{2})-([0-9]{2})\Z}
          \l_tmpa_tl \l__kanban_deadline_match_seq
        \int_set:Nn \l__kanban_reference_year_int
          {\seq_item:Nn \l__kanban_deadline_match_seq {2}}
        \int_set:Nn \l__kanban_reference_month_int
          {\seq_item:Nn \l__kanban_deadline_match_seq {3}}
        \int_set:Nn \l__kanban_reference_day_int
          {\seq_item:Nn \l__kanban_deadline_match_seq {4}}
      }
  }
\cs_new_protected:Npn \__kanban_classify_calendar_date:
  {
    \__kanban_load_reference_date:
    \int_set:Nn \l_tmpa_int
      {\l__kanban_deadline_year_int * 10000
        + \l__kanban_deadline_month_int * 100
        + \l__kanban_deadline_day_int}
    \int_set:Nn \l_tmpb_int
      {\l__kanban_reference_year_int * 10000
        + \l__kanban_reference_month_int * 100
        + \l__kanban_reference_day_int}
    \int_compare:nNnTF {\l_tmpa_int} < {\l_tmpb_int}
      {\booltrue{kb@overdue}}
      {
        \int_compare:nNnT {\l_tmpa_int} = {\l_tmpb_int}
          {\booltrue{kb@duenow}}
      }
  }
%    \end{macrocode}
%
% \paragraph{Project-month classification.}
% A project month begins on the configured start-day number, not necessarily on
% the first day of a calendar month. The classifier computes the nominal month
% offset from the project start, clamps the boundary day to the last valid day
% of shorter reference months, and subtracts one when the reference date has
% not yet reached that boundary. It then compares the resulting project-month
% number with the requested |M| value.
%
% The month offset is an ordinary LaTeX3 integer expression. |\int_set_eq:NN|
% temporarily places the reference year and month in the date component
% variables so the shared maximum-day function can be reused. The boundary is
% clamped through |\int_compare:nNnT| and another equality assignment.
% |\int_decr:N| adjusts the current project month before the final greater-than
% and equality tests set the same two semantic |etoolbox| Booleans used by
% calendar dates.
%    \begin{macrocode}
\cs_new_protected:Npn \__kanban_classify_project_month:
  {
    \__kanban_load_reference_date:
    \int_set:Nn \l_tmpa_int
      {(\l__kanban_reference_year_int
          - \l__kanban_project_start_year_int) * 12
        + (\l__kanban_reference_month_int
          - \l__kanban_project_start_month_int) + 1}
    \int_set_eq:NN \l__kanban_deadline_year_int
      \l__kanban_reference_year_int
    \int_set_eq:NN \l__kanban_deadline_month_int
      \l__kanban_reference_month_int
    \__kanban_set_max_day:
    \int_set_eq:NN \l__kanban_project_boundary_day_int
      \l__kanban_project_start_day_int
    \int_compare:nNnT
      {\l__kanban_project_boundary_day_int}
      > {\l__kanban_deadline_max_day_int}
      {\int_set_eq:NN \l__kanban_project_boundary_day_int
        \l__kanban_deadline_max_day_int}
    \int_compare:nNnT
      {\l__kanban_reference_day_int}
      < {\l__kanban_project_boundary_day_int}
      {\int_decr:N \l_tmpa_int}
    \int_compare:nNnTF {\l_tmpa_int} > {\l__kanban_project_month_int}
      {\booltrue{kb@overdue}}
      {
        \int_compare:nNnT
          {\l_tmpa_int} = {\l__kanban_project_month_int}
          {\booltrue{kb@duenow}}
      }
  }
%    \end{macrocode}
% \begin{macro}{\kb@evaloverdue}
% This internal entry point is responsible for reducing every deadline token
% list to two renderer-facing facts: overdue and due now. It resets both facts,
% tries the calendar grammar first, and otherwise distinguishes an invalid date
% from a project month and from an arbitrary neutral label. A project month is
% classified only when the board has a valid start-date snapshot. Card
% rendering therefore consumes only the two resulting Booleans and does not
% need to understand date syntax.
%
% |\boolfalse| clears the |etoolbox| flags at every call. The nested
% |\bool_if:NTF| tests branch on the LaTeX3 parser results. Valid calendar input
% delegates directly to the calendar classifier. Recognized but invalid syntax
% produces |\PackageWarning| and remains neutral. For a valid project month,
% the start date is parsed through the same component path and copied with
% |\int_set_eq:NN| before classification. |\MessageBreak| introduces a
% controlled line break in a long diagnostic without changing its meaning.
%    \begin{macrocode}
\cs_new_protected:Npn \kb@evaloverdue #1
  {
    \boolfalse{kb@overdue}
    \boolfalse{kb@duenow}
    \__kanban_parse_date:nn {#1} {\kanban@boarddateformat}
    \bool_if:NTF \l__kanban_valid_date_bool
      {\__kanban_classify_calendar_date:}
      {
        \bool_if:NTF \l__kanban_date_syntax_bool
          {\PackageWarning{kanban}
            {invalid~deadline~date~'#1';~displaying~a~neutral~label}}
          {
            \__kanban_parse_project_month:n {#1}
            \bool_if:NTF \l__kanban_valid_project_month_bool
              {
                \__kanban_parse_date:nn
                  {\kanban@thestart} {\kanban@boarddateformat}
                \bool_if:NT \l__kanban_valid_date_bool
                  {
                    \int_set_eq:NN \l__kanban_project_start_day_int
                      \l__kanban_deadline_day_int
                    \int_set_eq:NN \l__kanban_project_start_month_int
                      \l__kanban_deadline_month_int
                    \int_set_eq:NN \l__kanban_project_start_year_int
                      \l__kanban_deadline_year_int
                    \__kanban_classify_project_month:
                  }
              }
              {
                \bool_if:NT \l__kanban_project_month_syntax_bool
                  {\PackageWarning{kanban}
                    {invalid~project-month~deadline~'#1';
                      \MessageBreak neutral~style~used}}
              }
          }
      }
  }
%    \end{macrocode}
% \end{macro}
%
% \paragraph{Validated persistent and board-local date setters.}
% Start and reference dates enter the package through several configuration
% scopes, but all accepted non-empty values must be stored canonically. The
% persistent setters update defaults for future boards globally. The board
% setters update only the current environment. Empty values deliberately clear
% the corresponding setting, while invalid non-empty values retain the
% previously accepted value and issue the shared named warning.
%
% |\tl_if_blank:nTF| handles clearing separately from validation. Persistent
% values use the global TeX definitions |\gdef| and |\xdef|; board snapshots
% use the local |\def| and |\edef|. The |e| or |x| forms expand
% |\__kanban_canonical_date:| so the stored text is a self-contained canonical
% date rather than a reference to mutable parsing integers. |\msg_warning:nnnn|
% issues the message defined during bootstrap and supplies its two placeholders
% with the setting name and rejected value.
%    \begin{macrocode}
\cs_new_protected:Npn \kanban@setstartdate #1
  {
    \tl_if_blank:nTF {#1}
      {\gdef\kanban@startdate{}}
      {
        \__kanban_parse_date:nn {#1} {\kanban@dateformat}
        \bool_if:NTF \l__kanban_valid_date_bool
          {\xdef\kanban@startdate{\__kanban_canonical_date:}}
          {\msg_warning:nnnn {kanban}{invalid-date-setting}
            {startdate}{#1}}
      }
  }
\cs_new_protected:Npn \kanban@setboardstartdate #1
  {
    \tl_if_blank:nTF {#1}
      {\def\kanban@thestart{}}
      {
        \__kanban_parse_date:nn {#1} {\kanban@boarddateformat}
        \bool_if:NTF \l__kanban_valid_date_bool
          {\edef\kanban@thestart{\__kanban_canonical_date:}}
          {\msg_warning:nnnn {kanban}{invalid-date-setting}
            {board~startdate}{#1}}
      }
  }
\cs_new_protected:Npn \kanban@setreferencedate #1
  {
    \tl_if_blank:nTF {#1}
      {\gdef\kanban@referencedate{}}
      {
        \__kanban_parse_date:nn {#1} {\kanban@dateformat}
        \bool_if:NTF \l__kanban_valid_date_bool
          {\xdef\kanban@referencedate{\__kanban_canonical_date:}}
          {\msg_warning:nnnn {kanban}{invalid-date-setting}
            {reference-date}{#1}}
      }
  }
\cs_new_protected:Npn \kanban@setboardreferencedate #1
  {
    \tl_if_blank:nTF {#1}
      {\def\kanban@thereference{}}
      {
        \__kanban_parse_date:nn {#1} {\kanban@boarddateformat}
        \bool_if:NTF \l__kanban_valid_date_bool
          {\edef\kanban@thereference{\__kanban_canonical_date:}}
          {\msg_warning:nnnn {kanban}{invalid-date-setting}
            {board~reference-date}{#1}}
      }
  }
%    \end{macrocode}
%
% \paragraph{Safe dimension recognition.}
% The column-gap setters must reject malformed input before TeX attempts a
% dimension assignment, because an unrestricted assignment can consume later
% document tokens while searching for a unit. The validator accepts either a
% complete explicit dimension in a standard TeX unit or one existing dimension
% or skip register. It does not evaluate arbitrary macros or dimension
% expressions at this boundary.
%
% |\prg_new_protected_conditional:Npnn| defines a protected LaTeX3 conditional
% together with the requested |T|, |F|, and |TF| variants. The first helper uses
% |\token_if_dim_register:NTF| and |\token_if_skip_register:NTF| to recognize a
% single register token, and completes each path with |\prg_return_true:| or
% |\prg_return_false:|. The second helper trims the candidate and applies an
% anchored regular expression covering an optional sign, decimal syntax, and
% the standard TeX units. |(?:...)| groups without creating captures. If that
% grammar fails, |\tl_if_single_token:nTF| ensures that only one token is passed
% to the register conditional. The final two protected wrappers translate the
% LaTeX3 conditional and named message back to the LaTeX2e configuration layer.
% |\ExplSyntaxOff| then restores normal document syntax before the next group of
% traditional package definitions.
%    \begin{macrocode}
\prg_new_protected_conditional:Npnn
  \__kanban_dimension_register:N #1 {T,F,TF}
  {
    \token_if_dim_register:NTF #1
      {\prg_return_true:}
      {
        \token_if_skip_register:NTF #1
          {\prg_return_true:}
          {\prg_return_false:}
      }
  }
\prg_new_protected_conditional:Npnn
  \__kanban_dimension_valid:n #1 {T,F,TF}
  {
    \tl_set:Nx \l_tmpa_tl {#1}
    \tl_trim_spaces:N \l_tmpa_tl
    \exp_args:NnV \regex_match:nnTF
      {\A[+\-]?(?:[0-9]+(?:\.[0-9]*)?|\.[0-9]+)
        (?:pt|pc|in|bp|cm|mm|dd|cc|sp|ex|em)\Z}
      \l_tmpa_tl
      {\prg_return_true:}
      {
        \exp_args:NV \tl_if_single_token:nTF \l_tmpa_tl
          {
            \exp_args:NV \__kanban_dimension_register:NTF \l_tmpa_tl
              {\prg_return_true:}
              {\prg_return_false:}
          }
          {\prg_return_false:}
      }
  }
\cs_new_protected:Npn \kanban@validatedimensionTF #1#2#3
  {\__kanban_dimension_valid:nTF {#1}{#2}{#3}}
\cs_new_protected:Npn \kanban@dimensionwarning #1#2
  {\msg_warning:nnnn {kanban}{invalid-dimension}{#1}{#2}}
\ExplSyntaxOff
%    \end{macrocode}
%
% \subsection{Persistent configuration}
%
% Built-in palettes are registered through the public mechanism. Mode, date
% format, start date, reference date, column list, palette, semantic labels,
% and gap defaults then enter the dedicated |kanban/setup| LaTeX3 key family.
% Package options are accumulated and passed through the same transaction so
% package-level and document-level setup have identical validation and scope.
% The functions in this subsection distinguish persistent assignment from
% board-local assignment consistently: a package option or |\kbset| changes the
% source for later snapshots, whereas an environment key changes only the
% snapshot already being constructed.
%
% \paragraph{Built-in palette registration.}
% The package's default, muted, and monochrome palettes should obey exactly the
% same validation and registry rules as user-defined palettes. They are
% therefore registered through the canonical public command rather than by
% writing directly to internal property lists. The default palette is selected
% immediately after registration so a document with no configuration still
% produces the documented colors.
%
% These are ordinary calls to |\kanbandefinepalette|, whose two mandatory
% arguments provide the registry name and comma list. The percent sign after
% the opening brace suppresses an otherwise unwanted source-line space in the
% argument. |\kanban@selectpalette| then invokes the validated persistent
% selector described in the preceding subsection. Reusing public behavior in
% this way makes built-in and user palettes follow one code path.
%    \begin{macrocode}
\kanbandefinepalette{default}{%
  kanbanSlate,kanbanBlue,kanbanTeal,kanbanSage,
  kanbanAmber,kanbanClay,kanbanRose,kanbanPlum}
\kanbandefinepalette{muted}{kanbanMutedBlue,kanbanMutedGreen,
  kanbanMutedPurple,kanbanMutedBrown}
\kanbandefinepalette{monochrome}{black!85,black!75,black!65,black!55}
\kanban@selectpalette{default}
%    \end{macrocode}
%
% \paragraph{Mode selection.}
% Visual mode is a constrained textual choice. The persistent selector records
% a global |light| or |dark| default; the board selector records the same choice
% locally. Unknown values do not reset the previous selection.
%
% |\newcommand| defines a checked LaTeX2e command and |[1]| declares one
% argument. Both selectors first call the shared lowercase choice normalizer.
% |\ifdefstring| from |etoolbox| compares the replacement text of a macro with
% a literal string and accepts true and false branches. The global selector
% uses |\gdef|, while the board selector uses |\def| so its result disappears
% with the environment group. The final false branch emits a recoverable
% |\PackageWarning|.
%    \begin{macrocode}
\newcommand\kanban@selectmode[1]{%
  \kanban@normalizechoice{#1}%
  \ifdefstring{\kb@choice}{light}{\gdef\kanban@mode{light}}{%
    \ifdefstring{\kb@choice}{dark}{\gdef\kanban@mode{dark}}{%
      \PackageWarning{kanban}
        {unknown mode '#1'; keeping the current mode}}}}
\newcommand\kanban@selectboardmode[1]{%
  \kanban@normalizechoice{#1}%
  \ifdefstring{\kb@choice}{light}{\def\kanban@boardmode{light}}{%
    \ifdefstring{\kb@choice}{dark}{\def\kanban@boardmode{dark}}{%
      \PackageWarning{kanban}
        {unknown mode '#1'; keeping the current mode}}}}
%    \end{macrocode}
%
% \paragraph{Date-format canonicalization and selection.}
% Five canonical format names are supported, together with a small set of
% readable aliases. Canonicalization is separated from assignment so the same
% mapping serves both persistent and board-local selection. A recognized alias
% writes one of |us|, |english|, |german|, |iso|, or |ymd|; an unknown choice
% leaves the result empty and is diagnosed by the appropriate selector.
%
% |\def\kb@datechoice{}| initializes the result sentinel. Repeated
% |\ifdefstring| tests implement the alias table, and a successful test replaces
% the sentinel with its canonical value. |\ifdefempty| then distinguishes
% failure from success. |\global\let| assigns the persistent format macro the
% current meaning of the canonical result without retaining an intermediate
% expansion dependency; the local |\let| does the same for a board snapshot.
% |\MessageBreak| improves the line wrapping of the longer board warning.
%    \begin{macrocode}
\newcommand\kanban@canonicaldateformat[1]{%
  \kanban@normalizechoice{#1}%
  \def\kb@datechoice{}%
  \ifdefstring{\kb@choice}{us}{\def\kb@datechoice{us}}{}%
  \ifdefstring{\kb@choice}{american}{\def\kb@datechoice{us}}{}%
  \ifdefstring{\kb@choice}{english}{\def\kb@datechoice{english}}{}%
  \ifdefstring{\kb@choice}{uk}{\def\kb@datechoice{english}}{}%
  \ifdefstring{\kb@choice}{british}{\def\kb@datechoice{english}}{}%
  \ifdefstring{\kb@choice}{german}{\def\kb@datechoice{german}}{}%
  \ifdefstring{\kb@choice}{de}{\def\kb@datechoice{german}}{}%
  \ifdefstring{\kb@choice}{deutsch}{\def\kb@datechoice{german}}{}%
  \ifdefstring{\kb@choice}{iso}{\def\kb@datechoice{iso}}{}%
  \ifdefstring{\kb@choice}{ymd}{\def\kb@datechoice{ymd}}{}%
  \ifdefstring{\kb@choice}{year-first}{\def\kb@datechoice{ymd}}{}}
\newcommand\kanban@selectdateformat[1]{%
  \kanban@canonicaldateformat{#1}%
  \ifdefempty\kb@datechoice
    {\PackageWarning{kanban}
      {unknown date format '#1'; keeping the current date format}}
    {\global\let\kanban@dateformat\kb@datechoice}}
\newcommand\kanban@selectboarddateformat[1]{%
  \kanban@canonicaldateformat{#1}%
  \ifdefempty\kb@datechoice
    {\PackageWarning{kanban}
      {unknown board date format '#1';
        \MessageBreak keeping the current date format}}
    {\let\kanban@boarddateformat\kb@datechoice}}
%    \end{macrocode}
%
% \paragraph{Validated column-gap assignment.}
% A column gap must be a non-negative dimension. The persistent and board-local
% setters share the safe grammar established by the deadline subsystem, assign
% the accepted input to a candidate register, and compare the evaluated length
% with zero before committing it to the appropriate scope.
%
% |\kanban@validatedimensionTF| supplies success and failure branches.
% |\setlength| performs the actual LaTeX length assignment only after syntax
% validation. The primitive |\ifdim| compares dimensions; its true branch for a
% negative candidate issues the shared warning. A successful persistent update
% uses |\global| before the register assignment, while the board update remains
% local. The closing |\fi| terminates the primitive conditional.
%    \begin{macrocode}
\newcommand\kanban@setcolumnsep[1]{%
  \kanban@validatedimensionTF{#1}{%
    \setlength\kanban@candidatecolumnsep{#1}%
    \ifdim\kanban@candidatecolumnsep<0pt
      \kanban@dimensionwarning{columnsep}{#1}%
    \else
      \global\kanban@columnsep=\kanban@candidatecolumnsep
    \fi
  }{%
    \kanban@dimensionwarning{columnsep}{#1}}}
\newcommand\kanban@setboardcolumnsep[1]{%
  \kanban@validatedimensionTF{#1}{%
    \setlength\kanban@candidatecolumnsep{#1}%
    \ifdim\kanban@candidatecolumnsep<0pt
      \kanban@dimensionwarning{board columnsep}{#1}%
    \else
      \kanban@boardcolumnsep=\kanban@candidatecolumnsep
    \fi
  }{%
    \kanban@dimensionwarning{board columnsep}{#1}}}
%    \end{macrocode}
%
% \paragraph{Materialization of the visual mode.}
% The mode name is semantic state; rendering requires concrete text, surface,
% border, and default heading colors. The application function translates the
% active board mode into those values after all board options have been
% processed. Metadata chip colors are intentionally absent from this function
% because their documented light-theme styling remains stable in both modes.
%
% |\ifdefstring| selects the dark branch when the board-mode macro expands to
% |dark| and otherwise uses the light definitions. Each |\def| is local to the
% board environment. The stored values are named |xcolor| colors or mix
% expressions; later |tcolorbox| definitions expand them when drawing a card.
%    \begin{macrocode}
\newcommand\kanban@applymode{%
  \ifdefstring{\kanban@boardmode}{dark}{%
    \def\kanban@textcolor{white!88!black}%
    \def\kanban@cardcolor{kanbanInk}%
    \def\kanban@bordercolor{white!18!black}%
    \def\kanban@headingtext{white}%
  }{%
    \def\kanban@textcolor{kanbanInk}%
    \def\kanban@cardcolor{white}%
    \def\kanban@bordercolor{black!15}%
    \def\kanban@headingtext{white}%
  }}
%    \end{macrocode}
%
% \paragraph{Persistent and board-local semantic labels.}
% Priority and completion labels are customizable text, but an empty label
% would create an invisible chip with misleading spacing. The two setters
% therefore reject blank input and dynamically address either the persistent
% or board-local label family.
%
% |\ifblank| from |etoolbox| tests the argument without requiring it to be a
% predefined macro. |\csgdef| globally defines a control sequence whose name is
% assembled from character data, such as |kanban@highlabel|. |\csdef| performs
% the analogous local definition for names such as
% |kanban@boardhighlabel|. Dynamic |cs| forms keep four nearly identical label
% setters in one implementation while the key family supplies the validated
% label identifier.
%    \begin{macrocode}
\newcommand\kanban@setlabel[2]{%
  \ifblank{#2}
    {\PackageWarning{kanban}
      {label '#1' must not be empty; keeping the current value}}
    {\csgdef{kanban@#1label}{#2}}}
\newcommand\kanban@setboardlabel[2]{%
  \ifblank{#2}
    {\PackageWarning{kanban}
      {board label '#1' must not be empty; keeping the current value}}
    {\csdef{kanban@board#1label}{#2}}}
%    \end{macrocode}
%
% \paragraph{Nested label key families.}
% Label customization is expressed as a subordinate key list, for example
% |labels={high=URGENT,done=FINISHED}|. Separate persistent and board-local
% families direct identical key names to setters with different assignment
% scope. Each family also owns its unknown-key diagnostic, which identifies
% whether the error occurred in global setup or in a board override.
%
% |\ExplSyntaxOn| is re-enabled for the LaTeX3 key definitions.
% |\keys_define:nn| creates or extends a named family. A property written
% |.code:n| stores executable code receiving the key value as |#1|. The
% reserved |unknown| handler receives the misspelled key name through
% |\l_keys_key_str| and issues the named diagnostic with
% |\msg_warning:nnnn|. Defining nested families independently also prevents a
% label name such as |mode| from being confused with a top-level setup key.
%    \begin{macrocode}
\ExplSyntaxOn
\keys_define:nn {kanban/setup-labels}
  {
    high .code:n = {\kanban@setlabel{high}{#1}},
    medium .code:n = {\kanban@setlabel{medium}{#1}},
    low .code:n = {\kanban@setlabel{low}{#1}},
    done .code:n = {\kanban@setlabel{done}{#1}},
    unknown .code:n =
      {\msg_warning:nnnn {kanban}{unknown-key}
        {label}{\l_keys_key_str}}
  }
\keys_define:nn {kanban/board-labels}
  {
    high .code:n = {\kanban@setboardlabel{high}{#1}},
    medium .code:n = {\kanban@setboardlabel{medium}{#1}},
    low .code:n = {\kanban@setboardlabel{low}{#1}},
    done .code:n = {\kanban@setboardlabel{done}{#1}},
    unknown .code:n =
      {\msg_warning:nnnn {kanban}{unknown-key}
        {board~label}{\l_keys_key_str}}
  }
%    \end{macrocode}
%
% \paragraph{The persistent setup key family.}
% The main setup family is the single semantic entry point for package options
% and later calls to |\kanbansetup|. Immediate keys delegate to their validated
% selectors. Date values are staged rather than applied immediately, because a
% |date-format| key appearing later in the same comma list must take effect
% before an ambiguous start or reference date is interpreted. Alias keys keep
% legacy or concise spelling without duplicating implementation.
%
% In |l3keys|, |.code:n| executes an arbitrary handler with the supplied value.
% |.meta:n| rewrites one key as another key assignment: |dateformat| becomes
% |date-format| and |today| becomes |reference-date|. The date handlers use
% |\tl_set:Nn| to retain the literal value and |\bool_set_true:N| to record that
% the staged value must later be committed. The |labels| key calls
% |\keys_set:nn| recursively on the subordinate family. The family-level
% |unknown| handler supplies a uniform non-fatal warning for unsupported setup
% names.
%    \begin{macrocode}
\keys_define:nn {kanban/setup}
  {
    columns .code:n = {\kanban@selectcolumns{#1}},
    palette .code:n = {\kanban@selectpalette{#1}},
    date-format .code:n = {\kanban@selectdateformat{#1}},
    dateformat .meta:n = {date-format={#1}},
    startdate .code:n =
      {
        \tl_set:Nn \l__kanban_pending_startdate_tl {#1}
        \bool_set_true:N \l__kanban_pending_startdate_bool
      },
    reference-date .code:n =
      {
        \tl_set:Nn \l__kanban_pending_reference_date_tl {#1}
        \bool_set_true:N \l__kanban_pending_reference_date_bool
      },
    today .meta:n = {reference-date={#1}},
    mode .code:n = {\kanban@selectmode{#1}},
    columnsep .code:n = {\kanban@setcolumnsep{#1}},
    labels .code:n = {\keys_set:nn {kanban/setup-labels}{#1}},
    unknown .code:n =
      {\msg_warning:nnnn {kanban}{unknown-key}
        {setup}{\l_keys_key_str}}
  }
%    \end{macrocode}
%
% \paragraph{A deferred setup transaction.}
% Processing a setup list is a small transaction. It clears both staging flags,
% lets |l3keys| process the complete list in source order, and only then applies
% any staged dates under the final selected format. Consequently
% |startdate=07/08/2026,date-format=english| and the reversed key order have the
% same interpretation.
%
% |\keys_set:nn| parses and executes the comma-separated family input.
% |\bool_if:NT| executes each commit only when its staging flag is true.
% |\exp_args:NV| passes the stored value, rather than the token-list variable
% itself, to the relevant setter. The function is protected because key
% processing and assignment are not expandable operations. |\ExplSyntaxOff|
% restores the normal catcodes before the public LaTeX2e interface is declared.
%    \begin{macrocode}
\cs_new_protected:Npn \kanban@processsetup #1
  {
    \bool_set_false:N \l__kanban_pending_startdate_bool
    \bool_set_false:N \l__kanban_pending_reference_date_bool
    \keys_set:nn {kanban/setup}{#1}
    \bool_if:NT \l__kanban_pending_startdate_bool
      {\exp_args:NV \kanban@setstartdate
        \l__kanban_pending_startdate_tl}
    \bool_if:NT \l__kanban_pending_reference_date_bool
      {\exp_args:NV \kanban@setreferencedate
        \l__kanban_pending_reference_date_tl}
  }
\ExplSyntaxOff
%    \end{macrocode}
% \begin{macro}{\kanbansetup}
% \begin{macro}{\kbset}
% The canonical setup command exposes the transaction above with one mandatory
% key-list argument. The short alias follows the same collision-tolerant policy
% as the palette alias. This keeps the public namespace convenient without
% taking ownership of an existing command defined by the document class or
% another package.
%
% |\NewDocumentCommand| with signature |m| reads one mandatory argument and
% delegates it unchanged. The LaTeX2e test |\@ifundefined| chooses between
% installing the alias with |\let| and preserving the occupied name with a
% warning. Equivalence by |\let| ensures that |\kbset| and |\kanbansetup| have
% exactly the same argument parser and behavior.
%    \begin{macrocode}
\NewDocumentCommand\kanbansetup{m}{\kanban@processsetup{#1}}
\@ifundefined{kbset}
  {\let\kbset\kanbansetup}
  {\PackageWarning{kanban}
    {\string\kbset\space is already defined;
      \MessageBreak use \string\kanbansetup\space instead}}
%    \end{macrocode}
%
% \paragraph{Collection and application of package options.}
% LaTeX2e presents unknown package options one at a time through
% |\CurrentOption|. The package collects them into one token list and applies
% that list through the same deferred setup transaction only after
% |\ProcessOptions| has finished. Package options therefore share validation,
% aliases, key-order independence, and diagnostics with |\kanbansetup|.
%
% |\tl_new:N| allocates the global accumulator. |\tl_if_empty:NF| inserts a
% comma only after the first option, and |\tl_gput_right:Nn| appends globally.
% The expansion prefix |\exp_args:NNo| expands |\CurrentOption| once while
% retaining the destination variable. Application uses |\exp_args:NV| to pass
% the complete accumulated value to |\kanban@processsetup|. After leaving
% LaTeX3 syntax, |\DeclareOption*| installs the collector as the catch-all
% LaTeX2e option handler. |\ProcessOptions\relax| visits the requested options.
% The final call to |\kanban@applypackageoptions| commits the resulting setup
% list.
% Unknown package options consequently produce the same warning as unknown
% setup keys.
%    \begin{macrocode}
\ExplSyntaxOn
\tl_new:N \g__kanban_package_options_tl
\cs_new_protected:Npn \kanban@collectpackageoption
  {
    \tl_if_empty:NF \g__kanban_package_options_tl
      {\tl_gput_right:Nn \g__kanban_package_options_tl {,}}
    \exp_args:NNo \tl_gput_right:Nn
      \g__kanban_package_options_tl {\CurrentOption}
  }
\cs_new_protected:Npn \kanban@applypackageoptions
  {
    \tl_if_empty:NF \g__kanban_package_options_tl
      {\exp_args:NV \kanban@processsetup
        \g__kanban_package_options_tl}
  }
\ExplSyntaxOff
\DeclareOption*{\kanban@collectpackageoption}
\ProcessOptions\relax
\kanban@applypackageoptions
%    \end{macrocode}
% \end{macro}
% \end{macro}
%
% \subsection{Card components and renderers}
%
% Metadata elements are small inline boxes. Their semantic colors are stable
% across modes; only the card surface, card text, and card border respond to
% the light or dark board mode. Public, namespaced styles are appended after
% structural defaults. Fixed and breakable layouts reuse the same card renderer
% but provide different column containers. Rendering is deliberately delayed:
% task declarations store calls to these components, and the environment
% executes those calls only after it knows all column counts and dimensions.
%
% \paragraph{Structural and extension styles.}
% The first style registry distinguishes fixed package structure from supported
% user extension points. Empty public styles such as |kanban/card| may be
% extended with |\tcbset|, while private styles establish the common geometry
% of inline pills and the stable author and deadline colors.
%
% |\tcbset| defines |tcolorbox| keys globally. A |/.style| value expands to a
% reusable set of keys. |on line| makes a box participate in the current text
% line; |boxsep| and the side-specific dimensions control padding; |arc| rounds
% corners; |boxrule=0pt| removes the frame; |nobeforeafter| suppresses paragraph
% material; and |box align=base| aligns the box with surrounding text. The
% private |kanban@chip| and |kanban@due| styles inherit |kanban@pill| and add
% semantic foreground and background macros. Because the public extension key
% is applied last in each component definition below, a document may override
% presentation without replacing structural collection logic.
%    \begin{macrocode}
\tcbset{
  kanban/header/.style={},
  kanban/card/.style={},
  kanban/column-body/.style={},
  kanban/priority-chip/.style={},
  kanban/author-chip/.style={},
  kanban/deadline-chip/.style={},
  kanban/completed-chip/.style={},
  kanban@pill/.style={on line,boxsep=0pt,left=1.2mm,right=1.2mm,
    top=.4mm,bottom=.4mm,arc=.7mm,boxrule=0pt,nobeforeafter,
    box align=base},
  kanban@chip/.style={kanban@pill,colback=\kanban@authorback,
    coltext=\kanban@authortext},
  kanban@due/.style={kanban@pill,colback=\kanban@deadlineback,
    coltext=\kanban@deadlinetext}}
%    \end{macrocode}
%
% \paragraph{Priority and author chips.}
% A priority chip maps a normalized semantic value to the board-local label,
% semantic background, and accessible foreground. An empty or unknown value
% produces no chip. The author component has no semantic classification and
% therefore renders any non-blank author text with the stable author style.
%
% Both commands are checked LaTeX2e definitions made by |\newcommand|.
% |\ifblank| avoids creating empty inline boxes. The priority renderer initializes
% three temporary macros and fills them through |\ifstrequal| comparisons;
% |\ifdefempty| then guards the actual box. |\tcbox| draws a one-line
% |tcolorbox|, and its optional key list combines the structural pill, the
% selected colors, and the public extension style. |\bfseries| and
% |\scriptsize| are declarations scoped by the box content. The author renderer
% uses the inherited |kanban@chip| colors and applies |kanban/author-chip| last.
%    \begin{macrocode}
\newcommand\kanban@prioritypill[1]{\ifblank{#1}{}{%
  \def\kb@pl{}\def\kb@pc{}\def\kb@pt{}%
  \ifstrequal{#1}{high}{\def\kb@pl{\kanban@boardhighlabel}%
    \def\kb@pc{kanbanHigh}\def\kb@pt{\kanban@hightext}}{}%
  \ifstrequal{#1}{medium}{\def\kb@pl{\kanban@boardmediumlabel}%
    \def\kb@pc{kanbanMed}\def\kb@pt{\kanban@mediumtext}}{}%
  \ifstrequal{#1}{med}{\def\kb@pl{\kanban@boardmediumlabel}%
    \def\kb@pc{kanbanMed}\def\kb@pt{\kanban@mediumtext}}{}%
  \ifstrequal{#1}{low}{\def\kb@pl{\kanban@boardlowlabel}%
    \def\kb@pc{kanbanLow}\def\kb@pt{\kanban@lowtext}}{}%
  \ifdefempty\kb@pl{}{\tcbox[kanban@pill,
    colback=\kb@pc,colframe=\kb@pc,coltext=\kb@pt,
    kanban/priority-chip]{\bfseries\scriptsize\kb@pl}}}}
\newcommand\kanban@authorchip[1]{\ifblank{#1}{}{%
  \tcbox[kanban@chip,kanban/author-chip]{\scriptsize #1}}}
%    \end{macrocode}
%
% \paragraph{The deadline clock symbol.}
% Deadline chips include a compact clock drawn from primitives so the package
% does not depend on an icon font. Its dimensions are fixed in points and are
% small enough to align with script-sized metadata text.
%
% The inline |\tikz| command creates a self-contained picture. Its |baseline|
% option lowers the picture relative to the surrounding baseline. The first
% |\draw| command creates a circle, and the next two draw the clock hands; each
% uses the same thin |line width|. Coordinates and lengths are local to this
% picture, so the symbol does not alter the dimensions or style of later TikZ
% material.
%    \begin{macrocode}
\newcommand\kb@clock{\tikz[baseline=-.55ex]{%
  \draw[line width=.3pt] (0,0) circle (1.7pt);
  \draw[line width=.3pt] (0,0)--(0,1.1pt);
  \draw[line width=.3pt] (0,0)--(1pt,.2pt);}}
%    \end{macrocode}
%
% \paragraph{Deadline and completion chips.}
% Deadline rendering first classifies the raw value and then prepares its
% display notation. A manually |late| task or an automatically overdue date
% receives the high semantic color; a date equal to the reference date receives
% the medium color; future dates, project labels, and descriptive text retain
% the neutral deadline style. Completion is a separate Boolean flag rendered
% with the low semantic color. These colors remain identical in light and dark
% board modes.
%
% |\ifboolexpr| combines the string test for the stored late flag with the
% |etoolbox| overdue Boolean. Its |test| wrapper embeds |\ifstrequal|, while
% |or bool| reads the named Boolean. The nested |\ifbool| selects due-now versus
% neutral output. Each branch invokes |\tcbox| with the public deadline style
% last and includes the clock, spacing command |\,|, size, and optional bold
% declaration. |\kanban@completedchip| uses |\ifstrequal| to recognize the
% stored flag and otherwise expands to nothing.
%    \begin{macrocode}
\newcommand\kanban@deadlinechip[2]{\ifblank{#2}{}{%
  \kb@evaloverdue{#2}%
  \kanban@preparedeadline{#2}%
  \ifboolexpr{test {\ifstrequal{#1}{1}} or bool {kb@overdue}}
    {\tcbox[kanban@due,colback=kanbanHigh,coltext=\kanban@hightext,
        kanban/deadline-chip]
      {\kb@clock\,\scriptsize\bfseries \kb@displaydeadline}}
    {\ifbool{kb@duenow}
      {\tcbox[kanban@due,colback=kanbanMed,coltext=\kanban@mediumtext,
          kanban/deadline-chip]
        {\kb@clock\,\scriptsize\bfseries \kb@displaydeadline}}
      {\tcbox[kanban@due,kanban/deadline-chip]
        {\kb@clock\,\scriptsize \kb@displaydeadline}}}}}
\newcommand\kanban@completedchip[1]{\ifstrequal{#1}{1}
  {\tcbox[kanban@pill,colback=kanbanLow,colframe=kanbanLow,
      coltext=\kanban@lowtext,kanban/completed-chip]
    {\bfseries\scriptsize\kanban@boarddonelabel}}{}}
%    \end{macrocode}
%
% \paragraph{Priority-value validation.}
% Task options must store only priority values understood by the renderer. The
% setter normalizes case and surrounding spaces, accepts the four documented
% spellings, treats a blank choice as omission, and warns when non-blank input
% is unknown. The normalized value, rather than the original spelling, is
% stored with the task.
%
% The implementation reuses |\kanban@normalizechoice|. It initializes the
% task-priority macro |\kbtask@prio| with |\def|. Nested |\ifdefstring| tests select the accepted
% replacement. The final |\ifdefempty| distinguishes an intentional blank from
% a rejected choice before issuing |\PackageWarning|. Because all definitions
% occur while a task is being collected inside an environment group, the
% temporary value cannot leak to later tasks.
%    \begin{macrocode}
\newcommand\kanban@setpriority[1]{%
  \kanban@normalizechoice{#1}%
  \def\kbtask@prio{}%
  \ifdefstring{\kb@choice}{high}{\def\kbtask@prio{high}}{%
    \ifdefstring{\kb@choice}{medium}{\def\kbtask@prio{medium}}{%
      \ifdefstring{\kb@choice}{med}{\def\kbtask@prio{med}}{%
        \ifdefstring{\kb@choice}{low}{\def\kbtask@prio{low}}{%
          \ifdefempty\kb@choice{}{\PackageWarning{kanban}
            {unknown priority '#1'; omitting the priority chip}}}}}}}
%    \end{macrocode}
%
% \paragraph{Break-layout headers and shared task cards.}
% The page-breakable renderer needs a header box whose height is synchronized
% across its independent columns. All renderers need one full-width card whose
% surface and text follow the active mode and whose colored west rule identifies
% the target column. The shared card composer places task text first, then an
% optional priority/completion/author row, and finally an optional deadline row.
%
% |\newtcolorbox| defines an environment backed by |tcolorbox|. The header
% accepts two arguments for background and foreground colors. |enhanced| enables
% advanced drawing, |sharp corners| removes rounding, |fontupper| applies bold
% ragged-right text, and |equal height group| gives every header in the current
% board group the height of the tallest member. |valign=center| vertically
% centers shorter headings in that shared height. The card accepts its column
% color as one argument. |width=\linewidth| makes it fill the column;
% |borderline west| draws the semantic accent; the mode macros supply its
% surface, frame, and text; and |kanban/card| is applied last.
%
% |\kanban@rendercard| is a seven-argument deferred composer. |\notblank| from
% |etoolbox| executes metadata rows only when their concatenated fields contain
% material. |\par| begins a new row, |\vspace| separates it, and |\noindent|
% removes paragraph indentation. |\hfill| pushes the author chip to the right;
% a conditional |\hspace| separates adjacent priority and completion pills.
% The same composer is stored for both board layouts, which guarantees
% identical card semantics.
%    \begin{macrocode}
\newtcolorbox{kanban@header}[2]{enhanced,sharp corners,colback=#1,
  boxrule=0pt,coltext=#2,
  fontupper=\bfseries\raggedright,
  left=2mm,right=2mm,
  top=1.2mm,bottom=1.2mm,before skip=0pt,after skip=1.8mm,
  equal height group={KBH@\kanban@boardgroup},valign=center,
  kanban/header}
\newtcolorbox{kanban@card}[1]{enhanced,arc=.6mm,outer arc=.6mm,
  colback=\kanban@cardcolor,colframe=\kanban@bordercolor,
  coltext=\kanban@textcolor,boxrule=.4pt,width=\linewidth,
  borderline west={2.2pt}{0pt}{#1},left=2.4mm,right=2mm,top=1.6mm,
  bottom=1.6mm,before skip=0pt,after skip=1.8mm,
  fontupper=\scriptsize\raggedright,kanban/card}
\newcommand\kanban@rendercard[7]{%
  \begin{kanban@card}{#1}#2%
    \notblank{#3#4#7}{\par\vspace{1mm}\noindent
      \kanban@prioritypill{#3}%
      \notblank{#3}{\notblank{#7}{\hspace{1mm}}{}}{}%
      \kanban@completedchip{#7}\hfill\kanban@authorchip{#4}}{}%
    \notblank{#5}{\par\vspace{1mm}\noindent
      \kanban@deadlinechip{#6}{#5}}{}%
  \end{kanban@card}}
%    \end{macrocode}
%
% \paragraph{Fixed-board containers and renderer adapters.}
% The fixed renderer uses a two-row raster: the first row contains synchronized
% headings and the second contains equal-height transparent column bodies. The
% breakable renderer instead places each header and task list in a separate
% |paracol| stream. Three small adapter commands insert the dynamically stored
% column name, count, color, foreground, and deferred card list into the
% appropriate container.
%
% |kanban@fixedheader| is a |tcolorbox| environment without outer vertical
% skip; raster row equalization supplies its height. |kanban@fixedbody| hides
% the frame, uses |opacityback=0| for a transparent background, removes lateral
% padding, and draws only |borderline south| in a tint of the column color.
% In the adapters, |\csuse| from |etoolbox| constructs and expands a control
% sequence from character data, such as the color for numeric column |#1| or
% the count belonging to the active board group. The break adapter uses the
% equal-height |kanban@header| and emits the stored card list directly after it;
% its final |\par| closes the column material cleanly.
%    \begin{macrocode}
\newtcolorbox{kanban@fixedheader}[2]{enhanced,sharp corners,
  colback=#1,boxrule=0pt,coltext=#2,
  fontupper=\bfseries\raggedright,boxsep=0pt,left=2mm,right=2mm,
  top=1mm,bottom=1mm,before skip=0pt,after skip=0pt,valign=center,
  kanban/header}
\newtcolorbox{kanban@fixedbody}[1]{enhanced,sharp corners,frame hidden,
  colback=white,opacityback=0,boxrule=0pt,
  borderline south={.6pt}{0pt}{#1!55},boxsep=0pt,
  left=0pt,right=0pt,top=2mm,bottom=2mm,
  before skip=0pt,after skip=0pt,valign=top,kanban/column-body}
\newcommand\kanban@onehead[1]{%
  \begin{kanban@fixedheader}{\csuse{kanban@color@#1}}%
      {\csuse{kanban@headingtext@#1}}%
    \csuse{kanban@name@#1}\hfill
      \csuse{kanban@count@\kanban@boardgroup @#1}%
  \end{kanban@fixedheader}}
\newcommand\kanban@onecol[1]{%
  \begin{kanban@fixedbody}{\csuse{kanban@color@#1}}%
    \csuse{kanban@list@\kanban@boardgroup @#1}%
  \end{kanban@fixedbody}}
\newcommand\kanban@onecol@b[1]{%
  \begin{kanban@header}{\csuse{kanban@color@#1}}%
      {\csuse{kanban@headingtext@#1}}%
    \csuse{kanban@name@#1}\hfill
      \csuse{kanban@count@\kanban@boardgroup @#1}%
  \end{kanban@header}%
  \csuse{kanban@list@\kanban@boardgroup @#1}\par}
%    \end{macrocode}
%
% \subsection{Task collection}
%
% The task key family stores metadata without expanding the card text. A valid
% declaration appends one protected renderer call to the selected board-local
% list. Invalid or missing targets are diagnosed and never affect the counter.
% Collection and rendering are separate phases: this permits the environment
% to determine the tallest heading, the highest column, and the final column
% counts before any task card is typeset.
%
% \paragraph{Task-option keys.}
% The task family translates user-facing metadata into temporary task macros.
% Priority uses the validated semantic setter; author and deadline preserve
% their token lists; |due| aliases |deadline|; and |late| and |completed| are
% valueless flags. An unknown task option is ignored with a scoped warning.
%
% |\keys_define:nn| creates the |kanban/task| family. A |.code:n| handler sees
% the supplied value as |#1|. |.meta:n| redirects |due| through the canonical
% deadline handler. |.value_forbidden:n = true| tells |l3keys| that a Boolean
% flag must appear without |=...|, thereby distinguishing |completed| from a
% misleading form such as |completed=false|. The unknown handler uses
% |\l_keys_key_str| in the common named warning. The protected wrapper invokes
% |\keys_set:nn| and allows the surrounding LaTeX2e collector to remain outside
% LaTeX3 syntax.
%    \begin{macrocode}
\ExplSyntaxOn
\keys_define:nn {kanban/task}
  {
    prio .code:n = {\kanban@setpriority{#1}},
    author .code:n = {\def\kbtask@author{#1}},
    deadline .code:n = {\def\kbtask@deadline{#1}},
    due .meta:n = {deadline={#1}},
    late .code:n = {\def\kbtask@late{1}},
    late .value_forbidden:n = true,
    completed .code:n = {\def\kbtask@completed{1}},
    completed .value_forbidden:n = true,
    unknown .code:n =
      {\msg_warning:nnnn {kanban}{unknown-key}
        {task}{\l_keys_key_str}}
  }
\cs_new_protected:Npn \kanban@processtaskoptions #1
  {\keys_set:nn {kanban/task}{#1}}
\ExplSyntaxOff
%    \end{macrocode}
%
% \paragraph{Validation and deferred task collection.}
% The collector first gives every metadata field an empty value and processes
% the option list. It then requires an active board, a non-empty normalized
% target, and a target present in the board's identifier map. Each failed
% condition emits a specific error and ignores the task. A task inside a
% deliberately rejected nested board is silently discarded because the nesting
% error has already explained the cause; an ordinary task outside any board is
% diagnosed independently.
%
% |\ifbool| reads the |etoolbox| lifecycle flags. |\ifdefempty| tests the
% normalized column and identifier macros. |\PackageError| takes an error
% message and explanatory help; |\string| prints a command name without
% executing it, and |\MessageBreak| controls diagnostic layout. The nested
% branch structure ensures that neither counter nor list is modified until all
% target checks have succeeded.
%
% For a valid target, |\numgdef| globally defines the dynamic count macro to its
% current integer expression plus one. |\csxappto| globally appends expanded
% material to the dynamic task-list macro. Expansion is carefully selective:
% |\noexpand| retains the renderer command for the later rendering phase;
% |\csuse| resolves the already fixed column color; and |\unexpanded| protects
% the task body and all metadata from premature expansion. The
% |\expandafter| forms obtain the contents of temporary metadata macros before
% protecting those contents. The stored result is therefore a self-contained
% seven-argument call to |\kanban@rendercard| rather than a reference to task
% state that will be overwritten by the next declaration.
%    \begin{macrocode}
\newcommand\kanban@dotask[3][]{%
  \def\kbtask@prio{}\def\kbtask@author{}\def\kbtask@deadline{}%
  \def\kbtask@late{}\def\kbtask@completed{}%
  \kanban@processtaskoptions{#1}%
  \ifbool{kanban@active}
    {\kanban@normtaskcol{#2}%
     \ifdefempty\kbtask@column
       {\PackageError{kanban}
         {task has no assigned column; valid columns:
           \kanban@availablecolumns}
         {The first mandatory argument of
           \string\kanbantask\space must name a column from the
           enclosing board.\MessageBreak The task was ignored.}}
       {\ifdefempty\kbtask@columnid
       {\PackageError{kanban}
         {unknown kanban column '#2'; valid columns:
           \kanban@availablecolumns}
         {Column names are matched without regard to letter case.
           \MessageBreak
           The task was ignored because its normalized target is not in
           this board.}}
       {\expandafter\numgdef
           \csname kanban@count@\kanban@boardgroup
             @\kbtask@columnid\endcsname
           {\csname kanban@count@\kanban@boardgroup
             @\kbtask@columnid\endcsname+1}%
         \csxappto{kanban@list@\kanban@boardgroup
           @\kbtask@columnid}{%
           \noexpand\kanban@rendercard
           {\csuse{kanban@color@\kbtask@columnid}}{\unexpanded{#3}}%
           {\unexpanded\expandafter{\kbtask@prio}}%
           {\unexpanded\expandafter{\kbtask@author}}%
           {\unexpanded\expandafter{\kbtask@deadline}}%
           {\unexpanded\expandafter{\kbtask@late}}%
           {\unexpanded\expandafter{\kbtask@completed}}}%
       }}}
    {\ifbool{kanban@discardtasks}{}
      {\PackageError{kanban}
        {task declared outside a kanban environment}
        {Move the task declaration between \string\begin{kanban} and
          \string\end{kanban}. The task was ignored.}}}}
%    \end{macrocode}
% \begin{macro}{\kanbantask}
% \begin{macro}{\kbtask}
% The canonical public command supplies the compact task declaration syntax and
% delegates all key processing, validation, and collection to the internal
% implementation above. The concise command is installed by control-sequence
% equivalence only when its name is free; an existing definition is preserved
% with a warning.
%
% |\NewDocumentCommand| uses the signature |O{}mm|. |O{}| reads an optional
% bracketed key list and supplies an empty default when it is absent; the two
% |m| arguments read the target column and task body. The wrapper forwards the
% optional list in brackets because |\kanban@dotask| is a traditional LaTeX2e
% command with an optional argument. |\@ifundefined| performs the alias
% collision test, |\let| installs exact equivalence, and |\string| plus
% |\space| prints the affected command name legibly in the warning.
%    \begin{macrocode}
\NewDocumentCommand\kanbantask{O{}mm}{%
  \kanban@dotask[#1]{#2}{#3}}
\@ifundefined{kbtask}
  {\let\kbtask\kanbantask}
  {\PackageWarning{kanban}
    {\string\kbtask\space is already defined;
      \MessageBreak use \string\kanbantask\space instead}}
%    \end{macrocode}
% \end{macro}
% \end{macro}
%
% \subsection{Board lifecycle}
%
% Board keys update only the active board snapshot. Column setup allocates its
% isolated lists, prepares both rendering sequences, and validates the complete
% gap width before typesetting begins. The entry routine also rejects nesting
% before any outer state can be replaced. The environment itself supplies a TeX
% group, so assignments intentionally divided into persistent and local forms
% acquire the documented lifetime without a separate restoration procedure.
%
% \paragraph{Board column selection and board-option keys.}
% A blank board column value means ``use the persistent default'', while a
% non-blank value is normalized as a board-local override. The board key family
% then provides local counterparts to the persistent setup keys and adds the
% valueless |break| switch. Start and reference dates use the same deferred
% strategy as persistent setup so key order cannot alter ambiguous-date parsing.
%
% |\ifblank| selects default initialization or validated local selection.
% |\keys_define:nn| creates the |kanban/board| family. The |break| handler calls
% |\booltrue|, and |.value_forbidden:n| requires flag syntax without a value.
% |.code:n| handlers delegate to the board-local selectors; |.meta:n| provides
% the documented spelling aliases; and the date handlers stage values with
% |\tl_set:Nn| and |\bool_set_true:N|. The nested |labels| key invokes the
% board-label family. As in the other key families, the |unknown| handler uses
% |\l_keys_key_str| to identify the ignored option precisely.
%    \begin{macrocode}
\newcommand\kanban@setboardcolumns[1]{%
  \ifblank{#1}
    {\kanban@columnsboardinit}
    {\kanban@selectboardcolumns{#1}}}
\ExplSyntaxOn
\keys_define:nn {kanban/board}
  {
    break .code:n = {\booltrue{kanban@break}},
    break .value_forbidden:n = true,
    columns .code:n = {\kanban@setboardcolumns{#1}},
    date-format .code:n = {\kanban@selectboarddateformat{#1}},
    dateformat .meta:n = {date-format={#1}},
    startdate .code:n =
      {
        \tl_set:Nn \l__kanban_pending_startdate_tl {#1}
        \bool_set_true:N \l__kanban_pending_startdate_bool
      },
    reference-date .code:n =
      {
        \tl_set:Nn \l__kanban_pending_reference_date_tl {#1}
        \bool_set_true:N \l__kanban_pending_reference_date_bool
      },
    today .meta:n = {reference-date={#1}},
    palette .code:n = {\kanban@selectboardpalette{#1}},
    mode .code:n = {\kanban@selectboardmode{#1}},
    columnsep .code:n = {\kanban@setboardcolumnsep{#1}},
    labels .code:n = {\keys_set:nn {kanban/board-labels}{#1}},
    unknown .code:n =
      {\msg_warning:nnnn {kanban}{unknown-key}
        {board}{\l_keys_key_str}}
  }
%    \end{macrocode}
%
% \paragraph{A deferred board-option transaction.}
% Board-option processing mirrors persistent setup but commits staged dates to
% local snapshot setters. All immediate options, including a possible
% |date-format| override, are therefore resolved before the dates are parsed.
% The transaction runs after the persistent values have been copied into the
% board, so omitted keys naturally inherit their defaults.
%
% |\bool_set_false:N| resets the two staging sentinels, |\keys_set:nn| processes
% the complete option list, and |\bool_if:NT| performs only the requested
% commits. |\exp_args:NV| supplies each stored token-list value to its setter.
% The protected definition prevents assignment code from being expanded in a
% moving context. |\ExplSyntaxOff| returns to the LaTeX2e syntax used by the
% dynamic renderer-sequence construction below.
%    \begin{macrocode}
\cs_new_protected:Npn \kanban@processboardoptions #1
  {
    \bool_set_false:N \l__kanban_pending_startdate_bool
    \bool_set_false:N \l__kanban_pending_reference_date_bool
    \keys_set:nn {kanban/board}{#1}
    \bool_if:NT \l__kanban_pending_startdate_bool
      {\exp_args:NV \kanban@setboardstartdate
        \l__kanban_pending_startdate_tl}
    \bool_if:NT \l__kanban_pending_reference_date_bool
      {\exp_args:NV \kanban@setboardreferencedate
        \l__kanban_pending_reference_date_tl}
  }
\ExplSyntaxOff
%    \end{macrocode}
%
% \paragraph{Per-column state and renderer-sequence construction.}
% Each normalized column is represented internally by a numeric identifier.
% Setup stores its display name, initializes its board-qualified deferred task
% list and counter, assigns the cyclic palette background and foreground, and
% appends one adapter call to each renderer sequence. Board qualification is
% used for card lists and counts because those definitions are global during
% collection; simple color and name definitions are removed explicitly when
% the board closes.
%
% |\csdef| and |\csgdef| from |etoolbox| define dynamically named local and
% global control sequences. |\csedef| performs a local expanded definition, so
% the selected colors are fixed immediately. |\numdef| evaluates and stores the
% incremented number of columns. |\eappto| appends expanded material to the
% header, fixed-body, and break-layout sequences, while |\noexpand| preserves
% the adapter command for execution at rendering time. |\ifnumequal| treats the
% first break column specially; later columns prepend |\switchcolumn| so
% |paracol| receives one stream per board column.
%    \begin{macrocode}
\newcommand\kanban@setupcolumn[2]{%
  \csdef{kanban@name@#2}{#1}%
  \csgdef{kanban@list@\kanban@boardgroup @#2}{}%
  \csgdef{kanban@count@\kanban@boardgroup @#2}{0}%
  \kanban@palettecolor{#2}\csedef{kanban@color@#2}{\kb@thiscolor}%
  \csedef{kanban@headingtext@#2}{\kb@thisheadingtext}%
  \numdef\kanban@ncols{\kanban@ncols+1}%
  \eappto\kanban@headseq{\noexpand\kanban@onehead{#2}}%
  \eappto\kanban@colseq{\noexpand\kanban@onecol{#2}}%
  \ifnumequal{#2}{1}
    {\eappto\kanban@colseq@b{\noexpand\kanban@onecol@b{#2}}}
    {\eappto\kanban@colseq@b{\noexpand\switchcolumn
       \noexpand\kanban@onecol@b{#2}}}}
%    \end{macrocode}
%
% \paragraph{Whole-board geometry validation.}
% A syntactically valid non-negative gap may still be impossible for a board
% with many columns. Before fixed rendering, the validator computes the total
% width of the $N-1$ inter-column gaps and requires that total to be strictly
% smaller than |\linewidth|. An impossible value produces an actionable error
% and falls back to zero for this board, allowing the document to continue.
%
% The primitive |\ifnum| skips multiplication for a one-column board.
% Dimension-register assignment copies the accepted board gap to
% |\kanban@totalgaps|, and |\multiply| scales it by the integer expression
% |\numexpr\kanban@ncols-1\relax|. |\ifdim| compares the result with the
% available line width. Its empty true branch accepts valid geometry; the false
% branch calls |\PackageError| and locally assigns |0pt|. Each primitive
% conditional is terminated by its matching |\fi|.
%    \begin{macrocode}
\newcommand\kanban@validateboardgeometry{%
  \ifnum\kanban@ncols>1\relax
    \kanban@totalgaps=\kanban@boardcolumnsep
    \multiply\kanban@totalgaps by \numexpr\kanban@ncols-1\relax
    \ifdim\kanban@totalgaps<\linewidth\relax\else
      \PackageError{kanban}
        {column separation leaves no width for columns}
        {For N columns, (N-1) times columnsep must be smaller than
          \string\linewidth.\MessageBreak
          The current board will use columnsep=0pt.}%
      \kanban@boardcolumnsep=0pt
    \fi
  \fi}
%    \end{macrocode}
%
% \paragraph{Environment entry and state snapshot.}
% Opening a board must either reject unsupported nesting or create a complete,
% isolated snapshot before tasks are read. The successful path resets lifecycle
% flags, advances the unique board identifier, copies every persistent scalar
% and list default, processes local overrides, materializes the visual mode,
% clears renderer sequences, registers all columns, validates aggregate
% geometry, and suppresses source whitespace before the first task.
%
% The initial |\ifbool| detects an already active board before any outer state
% is overwritten. A rejected inner environment clears its render and active
% flags locally, sets |kanban@discardtasks|, and issues one |\PackageError|; its
% body can then be read safely without being attached to the outer board. On a
% valid entry, |\stepcounter| advances the LaTeX counter and |\arabic| supplies
% its printable value. |\edef| takes expanded scalar snapshots, while |\let|
% copies the current label meanings. Length-register assignment copies the
% persistent gap. The previously defined initialization and option functions
% then build the structured state. Finally, |\def| empties the three deferred
% sequences and resets the numeric column macro, and |\ignorespaces| removes
% indentation introduced by the source layout immediately after
% |\begin{kanban}|.
%    \begin{macrocode}
\newcommand\kanban@begin[1]{%
   \ifbool{kanban@active}
     {\boolfalse{kanban@render}%
      \boolfalse{kanban@active}%
      \booltrue{kanban@discardtasks}%
      \PackageError{kanban}
        {nested kanban environments are not supported}
        {Close the current board before opening another one.
          \MessageBreak The nested board was ignored.}}
     {\booltrue{kanban@render}%
   \boolfalse{kanban@break}%
   \boolfalse{kanban@discardtasks}%
   \booltrue{kanban@active}%
   \stepcounter{kanban@board}%
   \edef\kanban@boardgroup{\arabic{kanban@board}}%
   \edef\kanban@thestart{\kanban@startdate}%
   \edef\kanban@boarddateformat{\kanban@dateformat}%
   \edef\kanban@thereference{\kanban@referencedate}%
   \edef\kanban@boardmode{\kanban@mode}%
   \let\kanban@boardhighlabel\kanban@highlabel
   \let\kanban@boardmediumlabel\kanban@mediumlabel
   \let\kanban@boardlowlabel\kanban@lowlabel
   \let\kanban@boarddonelabel\kanban@donelabel
   \kanban@boardcolumnsep=\kanban@columnsep
   \kanban@paletteboardinit
   \kanban@columnsboardinit
   \kanban@processboardoptions{#1}%
   \kanban@applymode
   \def\kanban@ncols{0}\def\kanban@headseq{}%
   \def\kanban@colseq{}\def\kanban@colseq@b{}%
   \kanban@setupboardcolumns
   \kanban@validateboardgeometry
   \ignorespaces}}
%    \end{macrocode}
% \begin{macro}{\kanban}
% \begin{macro}{\endkanban}
% The public environment has one optional board-key list and no mandatory
% column argument; columns come from the current persistent default unless a
% board-local |columns| key overrides it. Closing the environment selects the
% page-breakable |paracol| renderer or the indivisible fixed raster. A rejected
% nested environment has its render flag cleared and therefore produces no
% board. Successfully rendered boards release their dynamically named state
% afterward.
%
% |\NewDocumentEnvironment| defines both the begin and end commands and checks
% that the environment name is new. Its signature |O{}| reads an optional key
% list with an empty default. The begin code delegates to |\kanban@begin|. In
% the end code, |\ifbool{kanban@render}| guards the complete rendering and
% cleanup path. |\par\medskip| leaves the task-declaration paragraph and
% separates the board from preceding text.
%
% In break mode, |\setlength| assigns the board gap to |paracol|'s standard
% |\columnsep|, |\begin{paracol}{\kanban@ncols}| opens the requested number of
% independent streams, and the prepared sequence supplies headers, cards, and
% |\switchcolumn| calls. In fixed mode, |\noindent| suppresses paragraph
% indentation and a full-line |minipage| makes the board indivisible. The
% |tcbraster| has one column per Kanban stage and two rows: all headers followed
% by all bodies. Its keys set the available width, board gap, zero row gap,
% equal height within each row, and top alignment. Executing
% |\kanban@headseq\kanban@colseq| fills the raster in row-major order.
% After either renderer, |\vspace{2mm}| provides documented following space and
% |\kanban@cleanupboard| undefines per-column state. The environment's ordinary
% TeX group then restores all remaining local snapshots automatically.
%    \begin{macrocode}
\NewDocumentEnvironment{kanban}{O{}}
  {\kanban@begin{#1}}
  {\ifbool{kanban@render}{\par\medskip
   \ifbool{kanban@break}
     {\setlength\columnsep{\kanban@boardcolumnsep}%
      \begin{paracol}{\kanban@ncols}\kanban@colseq@b\end{paracol}}
     {\noindent\begin{minipage}{\linewidth}%
      \begin{tcbraster}[
        raster columns=\kanban@ncols,
        raster rows=2,
        raster width=\linewidth,
        raster column skip=\kanban@boardcolumnsep,
        raster row skip=0pt,
        raster before skip=0pt,
        raster after skip=0pt,
        raster equal height=rows,
        raster valign=top]
        \kanban@headseq\kanban@colseq
      \end{tcbraster}%
      \end{minipage}}%
   \par\vspace{2mm}\kanban@cleanupboard}{}}
%</package>
%    \end{macrocode}
% \end{macro}
% \end{macro}
% \Finale
